% Options for packages loaded elsewhere
% Options for packages loaded elsewhere
\PassOptionsToPackage{unicode}{hyperref}
\PassOptionsToPackage{hyphens}{url}
%
\documentclass[
  11pt,
  letterpaper,
  oneside]{scrbook}
\usepackage{xcolor}
\usepackage[top=1in,bottom=1in,left=1.25in,right=1in]{geometry}
\usepackage{amsmath,amssymb}
\setcounter{secnumdepth}{5}
\usepackage{iftex}
\ifPDFTeX
  \usepackage[T1]{fontenc}
  \usepackage[utf8]{inputenc}
  \usepackage{textcomp} % provide euro and other symbols
\else % if luatex or xetex
  \usepackage{unicode-math} % this also loads fontspec
  \defaultfontfeatures{Scale=MatchLowercase}
  \defaultfontfeatures[\rmfamily]{Ligatures=TeX,Scale=1}
\fi
\usepackage{lmodern}
\ifPDFTeX\else
  % xetex/luatex font selection
\fi
% Use upquote if available, for straight quotes in verbatim environments
\IfFileExists{upquote.sty}{\usepackage{upquote}}{}
\IfFileExists{microtype.sty}{% use microtype if available
  \usepackage[]{microtype}
  \UseMicrotypeSet[protrusion]{basicmath} % disable protrusion for tt fonts
}{}
\usepackage{setspace}
\makeatletter
\@ifundefined{KOMAClassName}{% if non-KOMA class
  \IfFileExists{parskip.sty}{%
    \usepackage{parskip}
  }{% else
    \setlength{\parindent}{0pt}
    \setlength{\parskip}{6pt plus 2pt minus 1pt}}
}{% if KOMA class
  \KOMAoptions{parskip=half}}
\makeatother
% Make \paragraph and \subparagraph free-standing
\makeatletter
\ifx\paragraph\undefined\else
  \let\oldparagraph\paragraph
  \renewcommand{\paragraph}{
    \@ifstar
      \xxxParagraphStar
      \xxxParagraphNoStar
  }
  \newcommand{\xxxParagraphStar}[1]{\oldparagraph*{#1}\mbox{}}
  \newcommand{\xxxParagraphNoStar}[1]{\oldparagraph{#1}\mbox{}}
\fi
\ifx\subparagraph\undefined\else
  \let\oldsubparagraph\subparagraph
  \renewcommand{\subparagraph}{
    \@ifstar
      \xxxSubParagraphStar
      \xxxSubParagraphNoStar
  }
  \newcommand{\xxxSubParagraphStar}[1]{\oldsubparagraph*{#1}\mbox{}}
  \newcommand{\xxxSubParagraphNoStar}[1]{\oldsubparagraph{#1}\mbox{}}
\fi
\makeatother

\usepackage{color}
\usepackage{fancyvrb}
\newcommand{\VerbBar}{|}
\newcommand{\VERB}{\Verb[commandchars=\\\{\}]}
\DefineVerbatimEnvironment{Highlighting}{Verbatim}{commandchars=\\\{\}}
% Add ',fontsize=\small' for more characters per line
\usepackage{framed}
\definecolor{shadecolor}{RGB}{241,243,245}
\newenvironment{Shaded}{\begin{snugshade}}{\end{snugshade}}
\newcommand{\AlertTok}[1]{\textcolor[rgb]{0.68,0.00,0.00}{#1}}
\newcommand{\AnnotationTok}[1]{\textcolor[rgb]{0.37,0.37,0.37}{#1}}
\newcommand{\AttributeTok}[1]{\textcolor[rgb]{0.40,0.45,0.13}{#1}}
\newcommand{\BaseNTok}[1]{\textcolor[rgb]{0.68,0.00,0.00}{#1}}
\newcommand{\BuiltInTok}[1]{\textcolor[rgb]{0.00,0.23,0.31}{#1}}
\newcommand{\CharTok}[1]{\textcolor[rgb]{0.13,0.47,0.30}{#1}}
\newcommand{\CommentTok}[1]{\textcolor[rgb]{0.37,0.37,0.37}{#1}}
\newcommand{\CommentVarTok}[1]{\textcolor[rgb]{0.37,0.37,0.37}{\textit{#1}}}
\newcommand{\ConstantTok}[1]{\textcolor[rgb]{0.56,0.35,0.01}{#1}}
\newcommand{\ControlFlowTok}[1]{\textcolor[rgb]{0.00,0.23,0.31}{\textbf{#1}}}
\newcommand{\DataTypeTok}[1]{\textcolor[rgb]{0.68,0.00,0.00}{#1}}
\newcommand{\DecValTok}[1]{\textcolor[rgb]{0.68,0.00,0.00}{#1}}
\newcommand{\DocumentationTok}[1]{\textcolor[rgb]{0.37,0.37,0.37}{\textit{#1}}}
\newcommand{\ErrorTok}[1]{\textcolor[rgb]{0.68,0.00,0.00}{#1}}
\newcommand{\ExtensionTok}[1]{\textcolor[rgb]{0.00,0.23,0.31}{#1}}
\newcommand{\FloatTok}[1]{\textcolor[rgb]{0.68,0.00,0.00}{#1}}
\newcommand{\FunctionTok}[1]{\textcolor[rgb]{0.28,0.35,0.67}{#1}}
\newcommand{\ImportTok}[1]{\textcolor[rgb]{0.00,0.46,0.62}{#1}}
\newcommand{\InformationTok}[1]{\textcolor[rgb]{0.37,0.37,0.37}{#1}}
\newcommand{\KeywordTok}[1]{\textcolor[rgb]{0.00,0.23,0.31}{\textbf{#1}}}
\newcommand{\NormalTok}[1]{\textcolor[rgb]{0.00,0.23,0.31}{#1}}
\newcommand{\OperatorTok}[1]{\textcolor[rgb]{0.37,0.37,0.37}{#1}}
\newcommand{\OtherTok}[1]{\textcolor[rgb]{0.00,0.23,0.31}{#1}}
\newcommand{\PreprocessorTok}[1]{\textcolor[rgb]{0.68,0.00,0.00}{#1}}
\newcommand{\RegionMarkerTok}[1]{\textcolor[rgb]{0.00,0.23,0.31}{#1}}
\newcommand{\SpecialCharTok}[1]{\textcolor[rgb]{0.37,0.37,0.37}{#1}}
\newcommand{\SpecialStringTok}[1]{\textcolor[rgb]{0.13,0.47,0.30}{#1}}
\newcommand{\StringTok}[1]{\textcolor[rgb]{0.13,0.47,0.30}{#1}}
\newcommand{\VariableTok}[1]{\textcolor[rgb]{0.07,0.07,0.07}{#1}}
\newcommand{\VerbatimStringTok}[1]{\textcolor[rgb]{0.13,0.47,0.30}{#1}}
\newcommand{\WarningTok}[1]{\textcolor[rgb]{0.37,0.37,0.37}{\textit{#1}}}

\usepackage{longtable,booktabs,array}
\usepackage{calc} % for calculating minipage widths
% Correct order of tables after \paragraph or \subparagraph
\usepackage{etoolbox}
\makeatletter
\patchcmd\longtable{\par}{\if@noskipsec\mbox{}\fi\par}{}{}
\makeatother
% Allow footnotes in longtable head/foot
\IfFileExists{footnotehyper.sty}{\usepackage{footnotehyper}}{\usepackage{footnote}}
\makesavenoteenv{longtable}
\usepackage{graphicx}
\makeatletter
\newsavebox\pandoc@box
\newcommand*\pandocbounded[1]{% scales image to fit in text height/width
  \sbox\pandoc@box{#1}%
  \Gscale@div\@tempa{\textheight}{\dimexpr\ht\pandoc@box+\dp\pandoc@box\relax}%
  \Gscale@div\@tempb{\linewidth}{\wd\pandoc@box}%
  \ifdim\@tempb\p@<\@tempa\p@\let\@tempa\@tempb\fi% select the smaller of both
  \ifdim\@tempa\p@<\p@\scalebox{\@tempa}{\usebox\pandoc@box}%
  \else\usebox{\pandoc@box}%
  \fi%
}
% Set default figure placement to htbp
\def\fps@figure{htbp}
\makeatother





\setlength{\emergencystretch}{3em} % prevent overfull lines

\providecommand{\tightlist}{%
  \setlength{\itemsep}{0pt}\setlength{\parskip}{0pt}}



 


\usepackage{braket}
\usepackage{physics}
\usepackage{xcolor}
\definecolor{ketblue}{HTML}{1a6eb5}
\makeatletter
\@ifpackageloaded{bookmark}{}{\usepackage{bookmark}}
\makeatother
\makeatletter
\@ifpackageloaded{caption}{}{\usepackage{caption}}
\AtBeginDocument{%
\ifdefined\contentsname
  \renewcommand*\contentsname{Table of contents}
\else
  \newcommand\contentsname{Table of contents}
\fi
\ifdefined\listfigurename
  \renewcommand*\listfigurename{List of Figures}
\else
  \newcommand\listfigurename{List of Figures}
\fi
\ifdefined\listtablename
  \renewcommand*\listtablename{List of Tables}
\else
  \newcommand\listtablename{List of Tables}
\fi
\ifdefined\figurename
  \renewcommand*\figurename{Figure}
\else
  \newcommand\figurename{Figure}
\fi
\ifdefined\tablename
  \renewcommand*\tablename{Table}
\else
  \newcommand\tablename{Table}
\fi
}
\@ifpackageloaded{float}{}{\usepackage{float}}
\floatstyle{ruled}
\@ifundefined{c@chapter}{\newfloat{codelisting}{h}{lop}}{\newfloat{codelisting}{h}{lop}[chapter]}
\floatname{codelisting}{Listing}
\newcommand*\listoflistings{\listof{codelisting}{List of Listings}}
\makeatother
\makeatletter
\makeatother
\makeatletter
\@ifpackageloaded{caption}{}{\usepackage{caption}}
\@ifpackageloaded{subcaption}{}{\usepackage{subcaption}}
\makeatother
\usepackage{bookmark}
\IfFileExists{xurl.sty}{\usepackage{xurl}}{} % add URL line breaks if available
\urlstyle{same}
\hypersetup{
  pdftitle={Quantum Computing with ket},
  pdfauthor={Elliott Kirk},
  hidelinks,
  pdfcreator={LaTeX via pandoc}}


\title{Quantum Computing with ket}
\usepackage{etoolbox}
\makeatletter
\providecommand{\subtitle}[1]{% add subtitle to \maketitle
  \apptocmd{\@title}{\par {\large #1 \par}}{}{}
}
\makeatother
\subtitle{From classical bits to Shor's algorithm}
\author{Elliott Kirk}
\date{2026-03-13}
\begin{document}
\frontmatter
\maketitle

\renewcommand*\contentsname{Table of contents}
{
\setcounter{tocdepth}{2}
\tableofcontents
}

\setstretch{1.25}
\mainmatter
\bookmarksetup{startatroot}

\chapter*{Preface}\label{preface}
\addcontentsline{toc}{chapter}{Preface}

\markboth{Preface}{Preface}

This book teaches quantum computing by doing it.

Every concept is introduced alongside runnable code using
\href{https://github.com/dmvjs/ket}{\texttt{ket}} --- a TypeScript
quantum circuit simulator. You don't need a physics degree. You need
curiosity and a willingness to play with the examples.

\section*{How to read this book}\label{how-to-read-this-book}
\addcontentsline{toc}{section}{How to read this book}

\markright{How to read this book}

The book is structured in three parts:

\textbf{The Foundations} (Chapters 1--4) covers the vocabulary of
quantum computing: qubits, gates, superposition, and entanglement. These
chapters assume no physics background --- only familiarity with the idea
of a computer program.

\textbf{Algorithms} (Chapters 5--8) shows why quantum computers are
interesting: they can solve certain problems faster than any classical
computer. We'll build up to Shor's algorithm, which can factor large
numbers and is the reason quantum computers make cryptographers nervous.

\textbf{Reality} (Chapters 9--10) grounds everything in how quantum
hardware actually behaves. Real qubits are noisy. We'll see what that
means and how to work with it.

\section*{The interactive examples}\label{the-interactive-examples}
\addcontentsline{toc}{section}{The interactive examples}

\markright{The interactive examples}

In the HTML version of this book, code blocks run live in your browser.
Every visualization is produced by \texttt{ket} --- no screenshots, no
mock data. You can modify the code and re-run it.

In the PDF version, code is shown with its pre-rendered output.

\section*{A note on math}\label{a-note-on-math}
\addcontentsline{toc}{section}{A note on math}

\markright{A note on math}

Quantum computing has precise mathematical foundations. This book uses
that math when it clarifies rather than obscures. Key equations are
introduced slowly, with plain-English translations. If you've seen a
matrix before, you have enough math to read this book.

\begin{center}\rule{0.5\linewidth}{0.5pt}\end{center}

\emph{``If you think you understand quantum mechanics, you don't
understand quantum mechanics.''} --- Richard Feynman

(Feynman also said quantum computers might be the only way to simulate
nature efficiently. That idea still drives the field.)

\part{The Foundations}

\chapter{From Bits to Qubits}\label{sec-bits}

\begin{Shaded}
\begin{Highlighting}[]
\NormalTok{//| echo: false}
\NormalTok{ket = window.\_\_ket\_load}
\NormalTok{svg = (s) =\textgreater{} \{ const d = document.createElement("div"); d.innerHTML = s; return d; \}}
\end{Highlighting}
\end{Shaded}

A classical bit is a light switch --- off (0) or on (1). Every computer
you've ever used is built from billions of these, flipping furiously.

A qubit is stranger. Flip one and measure it: you get 0 or 1 with equal
probability. Every time, random. But apply the same flip \emph{twice}
and you always get back exactly where you started. No classical random
process can do that.

That's interference. It's what makes quantum computing interesting.

\section{Superposition}\label{superposition}

The \textbf{Hadamard gate} (\texttt{h}) takes a definite \(|0\rangle\)
and puts it into equal superposition:

\begin{Shaded}
\begin{Highlighting}[]
\NormalTok{c\_h = new ket.Circuit(1).h(0)}
\NormalTok{svg(c\_h.simulate(\{ shots: 10000 \}).toSVG())}
\end{Highlighting}
\end{Shaded}

50/50 --- each run varies slightly due to sampling, but the true
probability is exactly \(\frac{1}{2}\). Now apply it twice:

\begin{Shaded}
\begin{Highlighting}[]
\NormalTok{c\_hh = new ket.Circuit(1).h(0).h(0)}
\NormalTok{svg(c\_hh.simulate(\{ shots: 1000 \}).toSVG())}
\end{Highlighting}
\end{Shaded}

Always \texttt{\textbar{}0⟩}. The randomness cancels perfectly. This is
\textbf{interference} --- the defining feature that separates a qubit
from a coin flip.

\section{The math (briefly)}\label{the-math-briefly}

A qubit's state is written as:

\[|\psi\rangle = \alpha|0\rangle + \beta|1\rangle\]

\(\alpha\) and \(\beta\) are complex numbers called \textbf{amplitudes}.
The probability of measuring 0 is \(|\alpha|^2\), and of measuring 1 is
\(|\beta|^2\). They always sum to 1.

After \texttt{h(0)}, both amplitudes are \(\frac{1}{\sqrt{2}}\), so both
probabilities are \(\frac{1}{2}\).

You can read the amplitudes directly with \texttt{stateAsArray()}:

\begin{Shaded}
\begin{Highlighting}[]
\NormalTok{new ket.Circuit(1).h(0).stateAsArray()}
\end{Highlighting}
\end{Shaded}

Both amplitudes are exactly \(1/\sqrt{2} \approx 0.707\), with phase 0.
The \texttt{prob} column is \(|0.707|^2 = 0.5\). After a Z gate the
probabilities don't change --- but the phase does:

\begin{Shaded}
\begin{Highlighting}[]
\NormalTok{new ket.Circuit(1).h(0).z(0).stateAsArray()}
\end{Highlighting}
\end{Shaded}

Same probabilities, but the \(|1\rangle\) amplitude is now \(-0.707\)
(phase = \(\pi\)). Measurement can't see this --- but the next gate can.
Phase is invisible to the observer, but real to the computation.

The notation \(|\cdot\rangle\) (a \textbf{ket}) is standard physics
shorthand for a quantum state. This library is named after it.

\section{Measurement collapses the
state}\label{measurement-collapses-the-state}

Before you measure, the qubit is genuinely in superposition --- not
secretly one or the other. The moment you measure, it collapses to a
definite 0 or 1 and the amplitude information is gone.

This is why H·H works: both gates interact with the amplitudes
\emph{before} any measurement. Once you measure, the phase is lost.

To see it: a measurement between the two H gates collapses the qubit to
either \(|0\rangle\) or \(|1\rangle\). Then the second H acts on a
classical state --- and both \texttt{H\textbar{}0⟩} and
\texttt{H\textbar{}1⟩} give 50/50, not the original certainty:

\begin{Shaded}
\begin{Highlighting}[]
\NormalTok{svg(new ket.Circuit(1).h(0).simulate(\{ shots: 10000 \}).toSVG())}
\end{Highlighting}
\end{Shaded}

\begin{Shaded}
\begin{Highlighting}[]
\NormalTok{svg(new ket.Circuit(1).x(0).h(0).simulate(\{ shots: 10000 \}).toSVG())}
\end{Highlighting}
\end{Shaded}

Both are 50/50 --- the bars may differ by a percent or two from shot
noise, but the underlying probability is exactly \(\frac{1}{2}\) in both
cases. The second H has no memory of the first. Contrast with H·H, which
always returns \texttt{\textbar{}0⟩}.

\section{Bits vs.~qubits}\label{bits-vs.-qubits}

\begin{longtable}[]{@{}lll@{}}
\toprule\noalign{}
& Classical bit & Qubit \\
\midrule\noalign{}
\endhead
\bottomrule\noalign{}
\endlastfoot
State & Always 0 or 1 & Superposition until measured \\
Measurement effect & None & Collapses the state \\
Interference & No & Yes \\
\end{longtable}

A qubit doesn't store more information than a classical bit --- a
measurement still yields one bit. The advantage is what you can do with
amplitudes \emph{before} measurement: steer interference toward right
answers and away from wrong ones.

That's the rest of this book.

\section{What's next}\label{whats-next}

Chapter 2 introduces the \textbf{Bloch sphere} --- a geometric picture
of every possible single-qubit state --- and the full set of
single-qubit gates.

\chapter{Single-Qubit Gates}\label{sec-single-qubit}

\begin{Shaded}
\begin{Highlighting}[]
\NormalTok{//| echo: false}
\NormalTok{ket = window.\_\_ket\_load}
\NormalTok{svg = (s) =\textgreater{} \{ const d = document.createElement("div"); d.innerHTML = s; return d; \}}
\end{Highlighting}
\end{Shaded}

A qubit's state is a point on a sphere.

Every single-qubit gate is a rotation of that sphere. This isn't an
analogy --- it's the geometry. Understanding it makes every gate
intuitive.

\section{The Bloch sphere}\label{the-bloch-sphere}

The state \(|\psi\rangle = \alpha|0\rangle + \beta|1\rangle\) has two
real degrees of freedom (after accounting for normalisation and global
phase). Two angles --- \(\theta\) and \(\phi\) --- are enough to pin
down any single-qubit state:

\[|\psi\rangle = \cos\frac{\theta}{2}|0\rangle + e^{i\phi}\sin\frac{\theta}{2}|1\rangle\]

\(\theta\) runs from 0 (north pole, \(|0\rangle\)) to \(\pi\) (south
pole, \(|1\rangle\)). \(\phi\) is the angle around the equator. Together
they define a point on a unit sphere --- the \textbf{Bloch sphere}.

\begin{Shaded}
\begin{Highlighting}[]
\NormalTok{svg(new ket.Circuit(1).blochSphere(0))}
\end{Highlighting}
\end{Shaded}

\begin{Shaded}
\begin{Highlighting}[]
\NormalTok{svg(new ket.Circuit(1).x(0).blochSphere(0))}
\end{Highlighting}
\end{Shaded}

\begin{Shaded}
\begin{Highlighting}[]
\NormalTok{svg(new ket.Circuit(1).h(0).blochSphere(0))}
\end{Highlighting}
\end{Shaded}

From left to right: \(|0\rangle\) (north pole), \(|1\rangle\) (south
pole), \(H|0\rangle = |{+}\rangle\) (equator). The equator is equal
superposition. Every point in between is a valid qubit state.

\section{The Pauli gates --- axis
flips}\label{the-pauli-gates-axis-flips}

The three \textbf{Pauli gates} are 180° rotations around the three axes.

\textbf{X} rotates 180° around the X axis --- the quantum NOT gate,
swapping \(|0\rangle \leftrightarrow |1\rangle\):

\begin{Shaded}
\begin{Highlighting}[]
\NormalTok{svg(new ket.Circuit(1).x(0).blochSphere(0))}
\end{Highlighting}
\end{Shaded}

\begin{Shaded}
\begin{Highlighting}[]
\NormalTok{svg(new ket.Circuit(1).x(0).simulate(\{ shots: 1000 \}).toSVG())}
\end{Highlighting}
\end{Shaded}

\textbf{Z} rotates 180° around the Z axis. It leaves \(|0\rangle\)
unchanged but flips the sign of \(|1\rangle\):
\(Z|1\rangle = -|1\rangle\). This is a \textbf{phase flip} --- invisible
to single-qubit measurement, but it changes how the state interferes in
a larger circuit.

\textbf{Y} is X followed by Z (up to a global phase). Less common in
introductory circuits, but completes the Pauli set.

\section{H revisited --- a diagonal
rotation}\label{h-revisited-a-diagonal-rotation}

H is a 180° rotation around the axis halfway between X and Z. That's why
it maps poles to equator and back:

\begin{itemize}
\tightlist
\item
  \(H|0\rangle = |{+}\rangle\) → equator
\item
  \(H|{+}\rangle = |0\rangle\) → back to north pole
\item
  \(HH = I\) --- two 180° rotations around the same axis cancel
\end{itemize}

The last point explains Chapter 1's interference result. H is its own
inverse.

\section{Phase gates --- spinning around
Z}\label{phase-gates-spinning-around-z}

\textbf{S} and \textbf{T} are partial rotations around the Z axis: 90°
and 45° respectively. They don't change measurement probabilities ---
the state stays on the equator --- but they shift the phase \(\phi\),
which changes how the qubit interferes with others.

\begin{Shaded}
\begin{Highlighting}[]
\NormalTok{svg(new ket.Circuit(1).h(0).blochSphere(0))}
\end{Highlighting}
\end{Shaded}

\begin{Shaded}
\begin{Highlighting}[]
\NormalTok{svg(new ket.Circuit(1).h(0).t(0).blochSphere(0))}
\end{Highlighting}
\end{Shaded}

\begin{Shaded}
\begin{Highlighting}[]
\NormalTok{svg(new ket.Circuit(1).h(0).s(0).blochSphere(0))}
\end{Highlighting}
\end{Shaded}

Left to right: \texttt{H\textbar{}0⟩}, \texttt{T·H\textbar{}0⟩} (+45°),
\texttt{S·H\textbar{}0⟩} (+90°). Same latitude --- same measurement
probabilities --- different phase.

\section{Rotation gates --- arbitrary
angles}\label{rotation-gates-arbitrary-angles}

\texttt{rx}, \texttt{ry}, \texttt{rz} rotate by any angle around each
axis. The argument order is \texttt{(angle,\ qubit)}.

\begin{Shaded}
\begin{Highlighting}[]
\NormalTok{svg(new ket.Circuit(1).rx(Math.PI / 4, 0).blochSphere(0))}
\end{Highlighting}
\end{Shaded}

\begin{Shaded}
\begin{Highlighting}[]
\NormalTok{svg(new ket.Circuit(1).rx(Math.PI / 2, 0).blochSphere(0))}
\end{Highlighting}
\end{Shaded}

\begin{Shaded}
\begin{Highlighting}[]
\NormalTok{svg(new ket.Circuit(1).rx(Math.PI, 0).blochSphere(0))}
\end{Highlighting}
\end{Shaded}

\texttt{rx(π)} is exactly X. \texttt{rx(π/2)} is halfway --- a common
building block on real hardware, where every gate is ultimately
expressed as a combination of \texttt{rz} and \texttt{ry} rotations.

\section{Gate summary}\label{gate-summary}

\begin{longtable}[]{@{}llll@{}}
\toprule\noalign{}
Gate & Rotation & Effect on \(|0\rangle\) & Effect on \(|1\rangle\) \\
\midrule\noalign{}
\endhead
\bottomrule\noalign{}
\endlastfoot
X & 180° around X & → \(|1\rangle\) & → \(|0\rangle\) \\
Y & 180° around Y & → \(i|1\rangle\) & → \(-i|0\rangle\) \\
Z & 180° around Z & unchanged & → \(-|1\rangle\) \\
H & 180° around X+Z & → \(|{+}\rangle\) & → \(|{-}\rangle\) \\
S & 90° around Z & unchanged & → \(i|1\rangle\) \\
T & 45° around Z & unchanged & → \(e^{i\pi/4}|1\rangle\) \\
\end{longtable}

\section{What's next}\label{whats-next-1}

Chapter 3 adds a second qubit and introduces \textbf{entanglement} ---
the phenomenon that makes quantum computing qualitatively more powerful
than anything classical.

\chapter{Multiple Qubits}\label{sec-multi-qubit}

\begin{Shaded}
\begin{Highlighting}[]
\NormalTok{//| echo: false}
\NormalTok{ket = window.\_\_ket\_load}
\NormalTok{svg = (s) =\textgreater{} \{ const d = document.createElement("div"); d.innerHTML = s; return d; \}}
\NormalTok{pre = (s) =\textgreater{} \{ const d = document.createElement("pre"); d.textContent = s; return d; \}}
\end{Highlighting}
\end{Shaded}

One qubit has two basis states: \(|0\rangle\) and \(|1\rangle\).

Two qubits have four: \(|00\rangle\), \(|01\rangle\), \(|10\rangle\),
\(|11\rangle\).

\(n\) qubits have \(2^n\) basis states --- and a quantum state is a
superposition over \emph{all of them simultaneously}. This exponential
state space is where quantum computing gets its leverage.

\section{Two-qubit states}\label{two-qubit-states}

A two-qubit state is written:

\[|\psi\rangle = \alpha_{00}|00\rangle + \alpha_{01}|01\rangle + \alpha_{10}|10\rangle + \alpha_{11}|11\rangle\]

Four amplitudes, each complex, summing to 1 in squared magnitude.
Putting both qubits in equal superposition with H gives all four
outcomes equally likely:

\begin{Shaded}
\begin{Highlighting}[]
\NormalTok{c\_2h = new ket.Circuit(2).h(0).h(1)}
\NormalTok{svg(c\_2h.toSVG())}
\end{Highlighting}
\end{Shaded}

\begin{Shaded}
\begin{Highlighting}[]
\NormalTok{pre(c\_2h.draw())}
\end{Highlighting}
\end{Shaded}

\begin{Shaded}
\begin{Highlighting}[]
\NormalTok{svg(c\_2h.simulate(\{ shots: 2000 \}).toSVG())}
\end{Highlighting}
\end{Shaded}

Each of the four basis states gets amplitude \(\frac{1}{2}\),
probability \(\frac{1}{4}\). The histogram shows them equally
distributed. With 3 qubits it would be 8 outcomes. With 10, 1024.

\section{Individual qubits are still independent
here}\label{individual-qubits-are-still-independent-here}

When there's no two-qubit gate connecting them, each qubit evolves on
its own. We can see this on the Bloch sphere --- each qubit lives at its
own point, independently:

\begin{Shaded}
\begin{Highlighting}[]
\NormalTok{c\_ind = new ket.Circuit(2).h(0).x(1)}
\NormalTok{svg(c\_ind.toSVG())}
\end{Highlighting}
\end{Shaded}

Qubit 0 is at the equator (superposition). Qubit 1 is at the south pole
(\(|1\rangle\)).

\begin{Shaded}
\begin{Highlighting}[]
\NormalTok{svg(c\_ind.blochSphere(0))}
\end{Highlighting}
\end{Shaded}

\begin{Shaded}
\begin{Highlighting}[]
\NormalTok{svg(c\_ind.blochSphere(1))}
\end{Highlighting}
\end{Shaded}

Two different poles. The state is a product state --- it factors as
\(|{+}\rangle \otimes |1\rangle\). The qubits are completely
independent.

\section{The CNOT gate}\label{the-cnot-gate}

Single-qubit gates act on one qubit at a time. To create correlations
between qubits, you need a two-qubit gate.

The \textbf{CNOT} (controlled-NOT) is the workhorse:

\begin{itemize}
\tightlist
\item
  If the control is \(|0\rangle\), do nothing to the target
\item
  If the control is \(|1\rangle\), flip the target
\end{itemize}

\begin{Shaded}
\begin{Highlighting}[]
\NormalTok{svg(new ket.Circuit(2).cx(0, 1).toSVG())}
\end{Highlighting}
\end{Shaded}

\begin{Shaded}
\begin{Highlighting}[]
\NormalTok{pre(new ket.Circuit(2).cx(0, 1).draw())}
\end{Highlighting}
\end{Shaded}

The filled dot is the control; the ⊕ is the target. On classical inputs
it's just XOR. We can verify each row of the truth table by preparing
each input state and measuring:

\begin{Shaded}
\begin{Highlighting}[]
\NormalTok{cnot\_00 = new ket.Circuit(2).cx(0, 1)}
\NormalTok{svg(cnot\_00.simulate(\{ shots: 500 \}).toSVG())}
\end{Highlighting}
\end{Shaded}

\begin{Shaded}
\begin{Highlighting}[]
\NormalTok{cnot\_10 = new ket.Circuit(2).x(0).cx(0, 1)}
\NormalTok{svg(cnot\_10.simulate(\{ shots: 500 \}).toSVG())}
\end{Highlighting}
\end{Shaded}

\begin{Shaded}
\begin{Highlighting}[]
\NormalTok{cnot\_01 = new ket.Circuit(2).x(1).cx(0, 1)}
\NormalTok{svg(cnot\_01.simulate(\{ shots: 500 \}).toSVG())}
\end{Highlighting}
\end{Shaded}

\begin{Shaded}
\begin{Highlighting}[]
\NormalTok{cnot\_11 = new ket.Circuit(2).x(0).x(1).cx(0, 1)}
\NormalTok{svg(cnot\_11.simulate(\{ shots: 500 \}).toSVG())}
\end{Highlighting}
\end{Shaded}

Each histogram has a single bar --- the classical inputs map
deterministically to classical outputs. On superpositions, the CNOT does
something more interesting, which is the subject of Chapter 4.

\section{Building up circuits}\label{building-up-circuits}

Gates chain left to right. Each is a matrix acting on the joint state
vector. Here is a three-qubit circuit that puts each qubit through a
different single-qubit gate:

\begin{Shaded}
\begin{Highlighting}[]
\NormalTok{c\_3 = new ket.Circuit(3).h(0).x(1).rx(Math.PI / 3, 2)}
\NormalTok{svg(c\_3.toSVG())}
\end{Highlighting}
\end{Shaded}

\begin{Shaded}
\begin{Highlighting}[]
\NormalTok{pre(c\_3.draw())}
\end{Highlighting}
\end{Shaded}

\begin{Shaded}
\begin{Highlighting}[]
\NormalTok{svg(c\_3.simulate(\{ shots: 2000 \}).toSVG())}
\end{Highlighting}
\end{Shaded}

Qubit 0 is in superposition (H), qubit 1 is always \(|1\rangle\) (X),
qubit 2 has a partial rotation. The joint measurement reflects all
three: the second bit is always 1, the other two vary.

\section{Product states have zero
entanglement}\label{product-states-have-zero-entanglement}

When there are no two-qubit gates, the state is a product state and has
no entanglement across any bond. \texttt{bondEntropies()} gives the log₂
entanglement entropy across each consecutive pair of qubits:

\begin{Shaded}
\begin{Highlighting}[]
\NormalTok{c\_product = new ket.Circuit(4).h(0).x(1).h(2).rx(Math.PI/4, 3)}
\NormalTok{svg(c\_product.toSVG())}
\end{Highlighting}
\end{Shaded}

\begin{Shaded}
\begin{Highlighting}[]
\NormalTok{c\_product.bondEntropies()}
\end{Highlighting}
\end{Shaded}

All zeros. No gate has connected any qubit to any other. Each qubit is
an island. Chapter 4 shows what happens when CNOT introduces
correlations --- and how the entropy changes.

\section{Circuit depth and width}\label{circuit-depth-and-width}

Two dimensions describe the cost of a quantum circuit:

\begin{itemize}
\tightlist
\item
  \textbf{Width} --- number of qubits (space)
\item
  \textbf{Depth} --- longest path from input to output (time; determines
  exposure to noise)
\end{itemize}

Gates on separate qubits can run in parallel, reducing depth:

\begin{Shaded}
\begin{Highlighting}[]
\NormalTok{// depth 1: all three H gates run simultaneously}
\NormalTok{c\_parallel = new ket.Circuit(3).h(0).h(1).h(2)}
\NormalTok{svg(c\_parallel.toSVG())}
\end{Highlighting}
\end{Shaded}

\begin{Shaded}
\begin{Highlighting}[]
\NormalTok{// depth 3: same H gates but forced to run on one qubit serially}
\NormalTok{c\_serial = new ket.Circuit(1).h(0).h(0).h(0)}
\NormalTok{svg(c\_serial.toSVG())}
\end{Highlighting}
\end{Shaded}

The three-qubit version has depth 1 --- all three H gates execute
simultaneously. The single-qubit version has depth 3. Both apply three H
gates total, but the three-qubit circuit finishes three times faster on
hardware.

Real hardware runs gates in parallel where qubits don't interact.
Minimising depth is one of the key tasks in quantum compilation.

\section{What's next}\label{whats-next-2}

CNOT acting on a superposition creates something classical circuits
can't: \textbf{entanglement}. Chapter 4 explores what that means, how to
measure it, and why it's the engine behind every quantum speedup.

\chapter{Entanglement}\label{sec-entanglement}

\begin{Shaded}
\begin{Highlighting}[]
\NormalTok{//| echo: false}
\NormalTok{ket = window.\_\_ket\_load}
\NormalTok{svg = (s) =\textgreater{} \{ const d = document.createElement("div"); d.innerHTML = s; return d; \}}
\NormalTok{pre = (s) =\textgreater{} \{ const d = document.createElement("pre"); d.textContent = s; return d; \}}
\end{Highlighting}
\end{Shaded}

Apply H to qubit 0, then CNOT with qubit 0 as control and qubit 1 as
target:

\begin{Shaded}
\begin{Highlighting}[]
\NormalTok{bell = new ket.Circuit(2).h(0).cx(0, 1)}
\NormalTok{svg(bell.toSVG())}
\end{Highlighting}
\end{Shaded}

\begin{Shaded}
\begin{Highlighting}[]
\NormalTok{pre(bell.draw())}
\end{Highlighting}
\end{Shaded}

\begin{Shaded}
\begin{Highlighting}[]
\NormalTok{svg(bell.simulate(\{ shots: 2000 \}).toSVG())}
\end{Highlighting}
\end{Shaded}

Only \(|00\rangle\) and \(|11\rangle\). Never \(|01\rangle\) or
\(|10\rangle\).

The two qubits are \textbf{entangled}. Measure one and you instantly
know the other --- perfectly correlated, no matter how far apart they
are.

\section{What entanglement is}\label{what-entanglement-is}

Before measurement, neither qubit has a definite value. The joint state
is:

\[|\Phi^+\rangle = \frac{1}{\sqrt{2}}|00\rangle + \frac{1}{\sqrt{2}}|11\rangle\]

This is a \textbf{Bell state} --- the simplest maximally entangled
state. It cannot be written as a product of two independent qubit
states. There is no \(|\psi_0\rangle \otimes |\psi_1\rangle\) that
equals it.

That factorability is the test: if you can write the state as a product,
the qubits are independent. If you can't, they're entangled.

\section{Entanglement entropy}\label{entanglement-entropy}

There's a quantitative measure. \texttt{bondEntropies()} returns the
log₂ von Neumann entropy across each consecutive bond. A product state
scores 0. Maximal entanglement scores 1.

\begin{Shaded}
\begin{Highlighting}[]
\NormalTok{// Product state: no entanglement}
\NormalTok{product = new ket.Circuit(2).h(0).x(1)}
\NormalTok{product.bondEntropies()}
\end{Highlighting}
\end{Shaded}

\begin{Shaded}
\begin{Highlighting}[]
\NormalTok{// Bell state: maximum entanglement}
\NormalTok{bell.bondEntropies()}
\end{Highlighting}
\end{Shaded}

The product state scores exactly 0. The Bell state scores exactly 1 ---
one full bit of entanglement entropy across the bond between qubit 0 and
qubit 1. This is the physics. The histogram alone can't tell you this;
the entropy can.

\section{Looking at the Bloch sphere}\label{looking-at-the-bloch-sphere}

In a product state, each qubit has a well-defined Bloch vector:

\begin{Shaded}
\begin{Highlighting}[]
\NormalTok{svg(product.blochSphere(0))}
\end{Highlighting}
\end{Shaded}

\begin{Shaded}
\begin{Highlighting}[]
\NormalTok{svg(product.blochSphere(1))}
\end{Highlighting}
\end{Shaded}

Qubit 0 sits on the equator (superposition). Qubit 1 sits at the south
pole (\(|1\rangle\)). Each qubit has a pure state of its own.

Now look at the same qubits in the Bell state:

\begin{Shaded}
\begin{Highlighting}[]
\NormalTok{svg(bell.blochSphere(0))}
\end{Highlighting}
\end{Shaded}

\begin{Shaded}
\begin{Highlighting}[]
\NormalTok{svg(bell.blochSphere(1))}
\end{Highlighting}
\end{Shaded}

Both qubits sit at the \emph{origin} of the Bloch sphere --- not on the
surface, but at the centre. A point on the surface is a pure state. The
centre is a maximally mixed state. Each qubit, viewed alone, is
completely undetermined. The information lives only in the correlation
between them.

This is the geometric signature of entanglement: entangled qubits have
no individual Bloch vectors. The description of the pair cannot be
reduced to descriptions of two individuals.

\section{The four Bell states}\label{the-four-bell-states}

There are four maximally entangled two-qubit states, one for each
combination of relative phase and parity:

\[|\Phi^+\rangle = \tfrac{1}{\sqrt{2}}(|00\rangle + |11\rangle) \qquad |\Phi^-\rangle = \tfrac{1}{\sqrt{2}}(|00\rangle - |11\rangle)\]
\[|\Psi^+\rangle = \tfrac{1}{\sqrt{2}}(|01\rangle + |10\rangle) \qquad |\Psi^-\rangle = \tfrac{1}{\sqrt{2}}(|01\rangle - |10\rangle)\]

\begin{Shaded}
\begin{Highlighting}[]
\NormalTok{phi\_plus  = new ket.Circuit(2).h(0).cx(0,1)}
\NormalTok{phi\_minus = new ket.Circuit(2).h(0).cx(0,1).z(0)}
\NormalTok{psi\_plus  = new ket.Circuit(2).h(0).cx(0,1).x(1)}
\NormalTok{psi\_minus = new ket.Circuit(2).h(0).cx(0,1).x(1).z(0)}
\end{Highlighting}
\end{Shaded}

\begin{Shaded}
\begin{Highlighting}[]
\NormalTok{svg(phi\_plus.toSVG())}
\end{Highlighting}
\end{Shaded}

\begin{Shaded}
\begin{Highlighting}[]
\NormalTok{svg(phi\_plus.simulate(\{ shots: 2000 \}).toSVG())}
\end{Highlighting}
\end{Shaded}

\begin{Shaded}
\begin{Highlighting}[]
\NormalTok{svg(phi\_minus.toSVG())}
\end{Highlighting}
\end{Shaded}

\begin{Shaded}
\begin{Highlighting}[]
\NormalTok{svg(phi\_minus.simulate(\{ shots: 2000 \}).toSVG())}
\end{Highlighting}
\end{Shaded}

\begin{Shaded}
\begin{Highlighting}[]
\NormalTok{svg(psi\_plus.toSVG())}
\end{Highlighting}
\end{Shaded}

\begin{Shaded}
\begin{Highlighting}[]
\NormalTok{svg(psi\_plus.simulate(\{ shots: 2000 \}).toSVG())}
\end{Highlighting}
\end{Shaded}

\begin{Shaded}
\begin{Highlighting}[]
\NormalTok{svg(psi\_minus.toSVG())}
\end{Highlighting}
\end{Shaded}

\begin{Shaded}
\begin{Highlighting}[]
\NormalTok{svg(psi\_minus.simulate(\{ shots: 2000 \}).toSVG())}
\end{Highlighting}
\end{Shaded}

\(\Phi^+\) and \(\Phi^-\) show only \(|00\rangle\)/\(|11\rangle\).
\(\Psi^+\) and \(\Psi^-\) show only \(|01\rangle\)/\(|10\rangle\). The
\(\pm\) phase difference is invisible in the histograms --- it matters
when these states enter further gates. All four score entropy = 1:

\begin{Shaded}
\begin{Highlighting}[]
\NormalTok{[phi\_plus, phi\_minus, psi\_plus, psi\_minus].map(c =\textgreater{} c.bondEntropies())}
\end{Highlighting}
\end{Shaded}

\section{Why it's strange}\label{why-its-strange}

The correlations are stronger than anything classical physics allows.

Imagine two coins that always land the same way --- both heads or both
tails --- no matter how far apart they are. Classically, you'd explain
this by assuming they were preset: heads-heads or tails-tails, decided
before they separated. Hidden information.

Quantum mechanics rules this out. The Bell inequalities are a
mathematical test: correlations from hidden preset values can only be so
strong. Entangled particles violate them. The correlations are not
pre-agreed --- they're genuinely created at the moment of measurement.

This is what Einstein called ``spooky action at a distance.'' He thought
it meant quantum mechanics was incomplete. Experiments since the 1980s
have confirmed it isn't: entanglement is real, and the hidden-variable
explanation doesn't work.

\section{Three-qubit entanglement ---
GHZ}\label{three-qubit-entanglement-ghz}

The pattern extends. The \textbf{GHZ state} entangles three qubits:

\[|GHZ\rangle = \frac{1}{\sqrt{2}}|000\rangle + \frac{1}{\sqrt{2}}|111\rangle\]

\begin{Shaded}
\begin{Highlighting}[]
\NormalTok{ghz = new ket.Circuit(3).h(0).cx(0,1).cx(0,2)}
\NormalTok{svg(ghz.toSVG())}
\end{Highlighting}
\end{Shaded}

\begin{Shaded}
\begin{Highlighting}[]
\NormalTok{pre(ghz.draw())}
\end{Highlighting}
\end{Shaded}

\begin{Shaded}
\begin{Highlighting}[]
\NormalTok{svg(ghz.simulate(\{ shots: 2000 \}).toSVG())}
\end{Highlighting}
\end{Shaded}

Only \(|000\rangle\) and \(|111\rangle\). Measure any one qubit and the
other two are determined. The bond entropies tell the full story:

\begin{Shaded}
\begin{Highlighting}[]
\NormalTok{ghz.bondEntropies()}
\end{Highlighting}
\end{Shaded}

Both bonds carry entanglement. The three qubits are collectively
correlated across the whole register.

Now look at what the individual Bloch spheres show:

\begin{Shaded}
\begin{Highlighting}[]
\NormalTok{svg(ghz.blochSphere(0))}
\end{Highlighting}
\end{Shaded}

\begin{Shaded}
\begin{Highlighting}[]
\NormalTok{svg(ghz.blochSphere(1))}
\end{Highlighting}
\end{Shaded}

\begin{Shaded}
\begin{Highlighting}[]
\NormalTok{svg(ghz.blochSphere(2))}
\end{Highlighting}
\end{Shaded}

All three qubits are at the origin --- maximally mixed individually,
even though the joint state is pure. The information is entirely
non-local: it exists only in the three-way correlation.

\section{Entanglement as a resource}\label{entanglement-as-a-resource}

Entanglement isn't just philosophically interesting --- it's a
computational resource:

\begin{itemize}
\tightlist
\item
  \textbf{Quantum teleportation} --- transmit a qubit's full state using
  only classical bits, via a shared Bell pair
\item
  \textbf{Superdense coding} --- send two classical bits by physically
  transmitting one qubit
\item
  \textbf{Quantum error correction} --- detect and fix errors without
  measuring the data directly
\item
  \textbf{Every quantum speedup} --- Grover's, Shor's, and everything
  beyond --- depends on building and exploiting entanglement across many
  qubits
\end{itemize}

The exponential state space from Chapter 3 is only useful \emph{because}
entanglement lets you act on all of it coherently. Without entanglement,
a quantum computer is just a slow classical one.

\section{What's next}\label{whats-next-3}

Chapter 5 combines interference and entanglement into the first real
algorithm: Deutsch--Jozsa, a problem where a quantum computer beats any
classical computer in a single query.

\part{Algorithms}

\chapter{Quantum Algorithms}\label{sec-algorithms}

\begin{Shaded}
\begin{Highlighting}[]
\NormalTok{//| echo: false}
\NormalTok{ket = window.\_\_ket\_load}
\NormalTok{svg = (s) =\textgreater{} \{ const d = document.createElement("div"); d.innerHTML = s; return d; \}}
\end{Highlighting}
\end{Shaded}

The ingredients are in place: superposition, single-qubit gates,
multi-qubit states, entanglement. A quantum algorithm uses all of them
together --- amplifying the amplitudes of correct answers and cancelling
the rest.

The first clean example is \textbf{Deutsch--Jozsa}.

\section{The problem}\label{the-problem}

You're given a black-box function \(f : \{0,1\}^n \to \{0,1\}\). You're
promised it's either:

\begin{itemize}
\tightlist
\item
  \textbf{Constant} --- same output for every input (all 0s or all 1s)
\item
  \textbf{Balanced} --- exactly half the inputs map to 0, half to 1
\end{itemize}

Classically, worst case requires \(2^{n-1}+1\) queries --- you must
check more than half the inputs. Quantum mechanics does it in
\textbf{one query}.

\section{The circuit}\label{the-circuit}

The phase oracle encodes \(f\) as a phase flip \((-1)^{f(x)}\) on each
basis state \(|x\rangle\). The circuit is: H on all qubits → oracle → H
on all qubits → measure.

\begin{Shaded}
\begin{Highlighting}[]
\NormalTok{dj = (n, oracle) =\textgreater{} \{}
\NormalTok{  let c = new ket.Circuit(n)}
\NormalTok{  for (let q = 0; q \textless{} n; q++) c = c.h(q)}
\NormalTok{  c = oracle(c)}
\NormalTok{  for (let q = 0; q \textless{} n; q++) c = c.h(q)}
\NormalTok{  return c}
\NormalTok{\}}

\NormalTok{constantOracle = (c) =\textgreater{} c}
\NormalTok{balancedOracle = (c) =\textgreater{} c.cz(0, 1)}

\NormalTok{svg(dj(3, balancedOracle).toSVG())}
\end{Highlighting}
\end{Shaded}

\textbf{Constant oracle} --- all amplitude returns to \(|000\rangle\):

\begin{Shaded}
\begin{Highlighting}[]
\NormalTok{svg(dj(3, constantOracle).simulate(\{ shots: 500 \}).toSVG())}
\end{Highlighting}
\end{Shaded}

\textbf{Balanced oracle} --- \(|000\rangle\) amplitude cancels to zero:

\begin{Shaded}
\begin{Highlighting}[]
\NormalTok{svg(dj(3, balancedOracle).simulate(\{ shots: 500 \}).toSVG())}
\end{Highlighting}
\end{Shaded}

One circuit, one measurement, certain answer.

\section{Why it works}\label{why-it-works}

Apply H to all qubits: uniform superposition over all \(2^n\) inputs.
The oracle applies phase \((-1)^{f(x)}\) to each \(|x\rangle\).

\textbf{Constant}: every state gets the same phase. The second H
perfectly reverses the first --- all amplitude flows back to
\(|00\ldots0\rangle\).

\textbf{Balanced}: half the states get a phase flip. The symmetry is
broken. \(|00\ldots0\rangle\) accumulates destructive interference ---
its amplitude cancels exactly to zero.

This is interference doing useful work: one global property determined
by the coherent interaction of all \(2^n\) amplitudes simultaneously.

\section{Bernstein--Vazirani}\label{bernsteinvazirani}

A closely related algorithm finds a hidden \(n\)-bit string \(s\) with a
single query. The oracle computes \(f(x) = s \cdot x \pmod{2}\). The
circuit is identical to Deutsch--Jozsa --- the measurement directly
reads out \(s\):

\begin{Shaded}
\begin{Highlighting}[]
\NormalTok{// Hidden string s = 101 — phase oracle: Z on each qubit where s[i] = 1}
\NormalTok{bvOracle = (c) =\textgreater{} c.z(0).z(2)}
\NormalTok{svg(dj(3, bvOracle).simulate(\{ shots: 500 \}).toSVG())}
\end{Highlighting}
\end{Shaded}

The measurement gives \texttt{101} with certainty. Classically: \(n\)
queries. Quantum: one.

\section{What's next}\label{whats-next-4}

Deutsch--Jozsa solves a contrived problem. Chapter 6 covers
\textbf{Grover's algorithm} --- a quadratic speedup for unstructured
search, applicable to any problem expressible as ``find the input
satisfying this condition.''

\chapter{Grover's Search}\label{sec-grover}

\begin{Shaded}
\begin{Highlighting}[]
\NormalTok{//| echo: false}
\NormalTok{ket = window.\_\_ket\_load}
\NormalTok{svg = (s) =\textgreater{} \{ const d = document.createElement("div"); d.innerHTML = s; return d; \}}
\end{Highlighting}
\end{Shaded}

Search an unsorted list of \(N\) items for the one satisfying some
condition. Classically: \(O(N)\) queries worst case. Grover's algorithm:
\(O(\sqrt{N})\) queries.

That's a quadratic speedup, and it applies to \emph{any} problem where
you can write a function that recognises a correct answer.

\section{The idea}\label{the-idea}

Start in uniform superposition over all \(N = 2^n\) states. Each Grover
iteration does two things:

\begin{enumerate}
\def\labelenumi{\arabic{enumi}.}
\tightlist
\item
  \textbf{Oracle}: phase-flip the marked state --- multiply its
  amplitude by \(-1\)
\item
  \textbf{Diffusion}: reflect all amplitudes around their average ---
  amplifies the flipped state
\end{enumerate}

After \(\approx \frac{\pi}{4}\sqrt{N}\) iterations the marked amplitude
is near 1. Measure and you get the answer.

\section{Two qubits --- searching 4
items}\label{two-qubits-searching-4-items}

Searching for \(|11\rangle\). The oracle is a CZ (phase-flips
\(|11\rangle\)). One iteration suffices:

\begin{Shaded}
\begin{Highlighting}[]
\NormalTok{oracle\_11 = (c) =\textgreater{} c.cz(0, 1)}
\NormalTok{g2 = ket.grover(2, oracle\_11, 1)}
\NormalTok{svg(g2.toSVG())}
\end{Highlighting}
\end{Shaded}

\begin{Shaded}
\begin{Highlighting}[]
\NormalTok{svg(g2.simulate(\{ shots: 1000 \}).toSVG())}
\end{Highlighting}
\end{Shaded}

\(|11\rangle\) with certainty. One oracle call, one item in four.

\section{Three qubits --- searching 8
items}\label{three-qubits-searching-8-items}

Searching for \(|101\rangle\). The oracle flips the middle qubit to
convert \(|101\rangle \to |111\rangle\), applies a multi-controlled
phase flip, then undoes:

\begin{Shaded}
\begin{Highlighting}[]
\NormalTok{oracle\_101 = (c) =\textgreater{} c.x(1).h(2).ccx(0, 1, 2).h(2).x(1)}
\NormalTok{g3 = ket.grover(3, oracle\_101)   // optimal iterations chosen automatically}
\NormalTok{svg(g3.toSVG())}
\end{Highlighting}
\end{Shaded}

\begin{Shaded}
\begin{Highlighting}[]
\NormalTok{svg(g3.simulate(\{ shots: 1000 \}).toSVG())}
\end{Highlighting}
\end{Shaded}

\(|101\rangle\) dominates. Optimal iterations for \(n=3\):
\(\approx\frac{\pi}{4}\sqrt{8} \approx 2\).

\section{Over-iterating}\label{over-iterating}

Grover is a rotation, not a one-way boost. Iterate past the optimum and
the probability cycles back down:

\begin{Shaded}
\begin{Highlighting}[]
\NormalTok{g3\_over = ket.grover(3, oracle\_101, 5)}
\NormalTok{svg(g3\_over.simulate(\{ shots: 1000 \}).toSVG())}
\end{Highlighting}
\end{Shaded}

At 5 iterations --- past the optimum of 2 for \(n=3\) --- the
probability of \(|101\rangle\) has decayed from \textasciitilde93\% back
toward the noise floor. Stopping at the right point is part of the
algorithm.

\section{What Grover can and can't
do}\label{what-grover-can-and-cant-do}

Grover provides a \textbf{provable} \(\sqrt{N}\) speedup for
unstructured search. It speeds up brute-force components of many
classical problems: constraint satisfaction, collision finding, key
search.

What it can't do: provide an exponential speedup. Breaking a 128-bit
symmetric key classically costs \(2^{128}\) operations. With Grover:
\(2^{64}\) --- still astronomical. Symmetric cryptography needs longer
keys in a quantum world, not redesign.

Shor's algorithm, next, does achieve an exponential speedup --- but only
for structured problems like factoring.

\section{What's next}\label{whats-next-5}

Grover's speedup comes from amplitude amplification. Shor's exponential
speedup requires a more powerful subroutine: the \textbf{Quantum Fourier
Transform}. Chapter 7 builds it.

\chapter{The Quantum Fourier Transform}\label{sec-qft}

\begin{Shaded}
\begin{Highlighting}[]
\NormalTok{//| echo: false}
\NormalTok{ket = window.\_\_ket\_load}
\NormalTok{svg = (s) =\textgreater{} \{ const d = document.createElement("div"); d.innerHTML = s; return d; \}}
\end{Highlighting}
\end{Shaded}

The classical Fourier Transform decomposes a signal into its constituent
frequencies. The \textbf{Quantum Fourier Transform} does the same to
quantum amplitudes --- and it's the engine behind Shor's algorithm.

The key difference from the classical FFT: the QFT on \(n\) qubits uses
\(O(n^2)\) gates, while the best classical FFT on \(2^n\) data points
needs \(O(n \cdot 2^n)\) operations. Exponentially fewer gates.

The catch: you can't read out the transformed amplitudes directly ---
measurement collapses them. The QFT is useful as a subroutine that
restructures phases so subsequent gates can extract what you need.

\section{What it does}\label{what-it-does}

The QFT maps each basis state to a uniform superposition with phases
encoding the input:

\[\text{QFT}|j\rangle = \frac{1}{\sqrt{2^n}} \sum_{k=0}^{2^n-1} e^{2\pi i jk/2^n} |k\rangle\]

Every output state has equal amplitude --- measurement after a bare QFT
is uniform. The information is in the phases, which interfere with
subsequent gates.

\section{The circuit}\label{the-circuit-1}

\begin{Shaded}
\begin{Highlighting}[]
\NormalTok{qftCircuit = (n) =\textgreater{} \{}
\NormalTok{  let c = new ket.Circuit(n)}
\NormalTok{  ket.qft(c, Array.from(\{ length: n \}, (\_, i) =\textgreater{} i))}
\NormalTok{  return c}
\NormalTok{\}}
\NormalTok{svg(qftCircuit(4).toSVG())}
\end{Highlighting}
\end{Shaded}

The structure: H on each qubit, followed by controlled phase rotations
from subsequent qubits (angles \(\pi/2^k\) decreasing), then bit-order
reversal swaps. For \(n=4\): 10 gates. Classical FFT on 16 points: 32
multiplications.

\section{Phase estimation}\label{phase-estimation}

Phase estimation is the QFT's most important application. Given a
unitary \(U\) with eigenstate \(|\psi\rangle\) satisfying
\(U|\psi\rangle = e^{2\pi i\phi}|\psi\rangle\), it recovers \(\phi\) to
\(n\) bits of precision using \(n\) ancilla qubits and \(O(2^n)\)
applications of \(U\).

The S gate has eigenvalue \(e^{i\pi/2} = e^{2\pi i \cdot 1/4}\) on
\(|1\rangle\) --- phase \(\phi = 1/4\). With 4 precision qubits, the
encoded integer is \(\phi \cdot 2^4 = 4 = 0100_2\). In ket's qubit
ordering (qubit 0 is LSB), this registers as ancilla bits \texttt{0010}.
The histogram shows all 5 qubits --- the last bit is the target staying
in \(|1\rangle\):

\begin{Shaded}
\begin{Highlighting}[]
\NormalTok{qpe\_s = \{}
\NormalTok{  let c = new ket.Circuit(5).x(4)             // target qubit in eigenstate |1⟩}
\NormalTok{  for (let k = 0; k \textless{} 4; k++) c = c.h(k)     // superpose ancillas}
\NormalTok{  for (let k = 0; k \textless{} 4; k++)                 // controlled{-}S\^{}(2\^{}k)}
\NormalTok{    for (let p = 0; p \textless{} 2 ** k; p++) c = c.cs(k, 4)}
\NormalTok{  return ket.applyIqft(c, 4, 0)               // inverse QFT}
\NormalTok{\}}
\NormalTok{svg(qpe\_s.toSVG())}
\end{Highlighting}
\end{Shaded}

\begin{Shaded}
\begin{Highlighting}[]
\NormalTok{svg(qpe\_s.simulate(\{ shots: 500 \}).toSVG())}
\end{Highlighting}
\end{Shaded}

\texttt{00101} with certainty --- the 4 ancilla bits read \texttt{0010},
meaning qubit 2 is set. In LSB-first encoding that's integer 4, and
\(4/16 = 1/4\). Phase recovered exactly. One run, four bits of
precision.

\section{Connection to Shor's}\label{connection-to-shors}

Shor's algorithm reduces factoring to period finding: given
\(f(x) = a^x \bmod N\), find its period \(r\). Period finding is phase
estimation on the unitary \(U_a|x\rangle = |ax \bmod N\rangle\) --- its
eigenvalues encode \(1/r\).

Building the circuit for \(U_a\) is the hard part (the Beauregard
circuit). Once you have it, QPE extracts the period, and from the
period, classical arithmetic gives the factors. Chapter 8 runs the whole
thing.

\chapter{Shor's Algorithm}\label{sec-shor}

\begin{Shaded}
\begin{Highlighting}[]
\NormalTok{//| echo: false}
\NormalTok{ket = window.\_\_ket\_load}
\NormalTok{svg = (s) =\textgreater{} \{ const d = document.createElement("div"); d.innerHTML = s; return d; \}}
\NormalTok{pre = (s) =\textgreater{} \{ const d = document.createElement("pre"); d.textContent = s; return d; \}}
\end{Highlighting}
\end{Shaded}

In 1994, Peter Shor showed that a quantum computer can factor large
integers in polynomial time. This broke the assumption underpinning RSA
encryption and started the modern era of quantum computing research.

\section{Why factoring matters}\label{why-factoring-matters}

RSA security rests on a simple asymmetry: multiplying two large primes
is easy, factoring their product is hard. The best known classical
algorithm runs in roughly \(e^{O(n^{1/3})}\) time for an \(n\)-bit
number. Shor's algorithm runs in \(O(n^3)\) --- polynomial. A large
enough quantum computer breaks RSA.

\section{The reduction}\label{the-reduction}

Shor doesn't factor directly. It reduces factoring to \textbf{period
finding}:

\begin{enumerate}
\def\labelenumi{\arabic{enumi}.}
\tightlist
\item
  Pick random \(a < N\) with \(\gcd(a, N) = 1\)
\item
  Find period \(r\) of \(f(x) = a^x \bmod N\)
\item
  Then \(\gcd(a^{r/2} \pm 1,\, N)\) gives the factors (if \(r\) is even
  and \(a^{r/2} \not\equiv -1\))
\end{enumerate}

Steps 1 and 3 are classical. Step 2 is quantum --- period finding via
phase estimation on the unitary \(U_a|x\rangle = |ax \bmod N\rangle\).

\section{Running it}\label{running-it}

\texttt{shorBeauregard} implements the full gate-decomposed Beauregard
circuit:

\begin{Shaded}
\begin{Highlighting}[]
\NormalTok{r15 = ket.shorBeauregard(15n, 7n)}
\NormalTok{pre(\textasciigrave{}factors = $\{r15.factors.map(String).join(\textquotesingle{} \textbackslash{}xd7 \textquotesingle{})\}  (base $\{r15.a\}, $\{r15.qubits\} qubits)\textasciigrave{})}
\end{Highlighting}
\end{Shaded}

The algorithm selects the most efficient base internally. Here it
factors 15 into 3 and 5 using 19 qubits.

\begin{Shaded}
\begin{Highlighting}[]
\NormalTok{r21 = ket.shorBeauregard(21n, 2n)}
\NormalTok{pre(\textasciigrave{}21 = $\{r21.factors.map(String).join(\textquotesingle{} \textbackslash{}xd7 \textquotesingle{})\}  ($\{r21.qubits\} qubits)\textasciigrave{})}
\end{Highlighting}
\end{Shaded}

\begin{Shaded}
\begin{Highlighting}[]
\NormalTok{r35 = ket.shorBeauregard(35n, 3n)}
\NormalTok{pre(\textasciigrave{}35 = $\{r35.factors.map(String).join(\textquotesingle{} \textbackslash{}xd7 \textquotesingle{})\}  ($\{r35.qubits\} qubits)\textasciigrave{})}
\end{Highlighting}
\end{Shaded}

\section{The qubit cost}\label{the-qubit-cost}

The Beauregard circuit for factoring \(N\) requires \(O(\log N)\)
qubits. For \(N=15\): 19 qubits. For a 2048-bit RSA key: roughly 4000
logical qubits, plus thousands more for error correction. Current
hardware has hundreds of noisy physical qubits.

The algorithm works. The hardware isn't there yet.

\section{Simulation limits}\label{simulation-limits}

Simulation is feasible only for small \(N\) --- the circuit's state
space grows with \(\log N\), but intermediate entanglement makes
classical simulation exponentially expensive:

\begin{longtable}[]{@{}lll@{}}
\toprule\noalign{}
N & Qubits & Approx. simulation time \\
\midrule\noalign{}
\endhead
\bottomrule\noalign{}
\endlastfoot
15 & 19 & \textless{} 1s \\
21 & 22 & \textasciitilde3s \\
35 & 25 & \textasciitilde30s \\
77 & 29 & \textasciitilde15 min \\
\end{longtable}

This is the honest limit of classical simulation. To factor a
cryptographically relevant number you need actual quantum hardware ---
which is the entire point.

\section{Try it --- QPE factoring}\label{try-it-qpe-factoring}

This interactive demo runs the full quantum period-finding circuit in
the browser. Select any number from the list --- including 4-digit
semiprimes --- and watch the QPE counting register reveal the period and
the factors.

The approach: build a permutation unitary
\(U_a|x\rangle = |ax \bmod N\rangle\) directly, embed it as a controlled
gate, run QPE with 3 precision bits, then extract the period via
continued fractions and compute \(\gcd(a^{r/2} \pm 1, N)\). Numbers
marked ⚡ use 12 work qubits and take 5--10 seconds.

\begin{Shaded}
\begin{Highlighting}[]
\NormalTok{//| echo: false}
\NormalTok{shor\_gcd = (a, b) =\textgreater{} \{ while (b) [a, b] = [b, a \% b]; return a \}}

\NormalTok{shor\_mulMatrix = (mul, N, n) =\textgreater{} \{}
\NormalTok{  const dim = 1 \textless{}\textless{} n}
\NormalTok{  const M = Array.from(\{ length: dim \}, () =\textgreater{} new Array(dim).fill(0))}
\NormalTok{  for (let x = 0; x \textless{} dim; x++) M[x \textless{} N ? x * mul \% N : x][x] = 1}
\NormalTok{  return M}
\NormalTok{\}}

\NormalTok{shor\_controlled = (U, n) =\textgreater{} \{}
\NormalTok{  const dim = 1 \textless{}\textless{} n}
\NormalTok{  const C = Array.from(\{ length: 2 * dim \}, (\_, i) =\textgreater{} \{}
\NormalTok{    const row = new Array(2 * dim).fill(0)}
\NormalTok{    if (i \textless{} dim) row[i] = 1}
\NormalTok{    return row}
\NormalTok{  \})}
\NormalTok{  for (let i = 0; i \textless{} dim; i++)}
\NormalTok{    for (let j = 0; j \textless{} dim; j++)}
\NormalTok{      C[dim + i][dim + j] = U[i][j]}
\NormalTok{  return C}
\NormalTok{\}}

\NormalTok{SHOR\_TARGETS = [}
\NormalTok{  \{ N: 15,   a: 2,   aSq: 4,    n: 4,  label: "15 = 3 × 5" \},}
\NormalTok{  \{ N: 21,   a: 8,   aSq: 1,    n: 5,  label: "21 = 7 × 3" \},}
\NormalTok{  \{ N: 33,   a: 10,  aSq: 1,    n: 6,  label: "33 = 3 × 11" \},}
\NormalTok{  \{ N: 35,   a: 6,   aSq: 1,    n: 6,  label: "35 = 5 × 7" \},}
\NormalTok{  \{ N: 77,   a: 34,  aSq: 1,    n: 7,  label: "77 = 7 × 11" \},}
\NormalTok{  \{ N: 1073, a: 813, aSq: 1,    n: 11, label: "1073 = 29 × 37" \},}
\NormalTok{  \{ N: 1517, a: 778, aSq: 1,    n: 11, label: "1517 = 37 × 41" \},}
\NormalTok{  \{ N: 1591, a: 1074,aSq: 1,    n: 11, label: "1591 = 37 × 43" \},}
\NormalTok{  \{ N: 1763, a: 42,  aSq: 1,    n: 11, label: "1763 = 41 × 43" \},}
\NormalTok{  \{ N: 1927, a: 565, aSq: 1270, n: 11, label: "1927 = 41 × 47" \},}
\NormalTok{  \{ N: 2021, a: 988, aSq: 1,    n: 11, label: "2021 = 43 × 47" \},}
\NormalTok{  \{ N: 2033, a: 322, aSq: 1,    n: 11, label: "2033 = 19 × 107" \},}
\NormalTok{  \{ N: 2047, a: 390, aSq: 622,  n: 11, label: "2047 = 23 × 89" \},}
\NormalTok{  \{ N: 2491, a: 847, aSq: 1,    n: 12, label: "2491 = 47 × 53 ⚡" \},}
\NormalTok{  \{ N: 2701, a: 2554,aSq: 1,    n: 12, label: "2701 = 37 × 73 ⚡" \},}
\NormalTok{]}
\end{Highlighting}
\end{Shaded}

\begin{Shaded}
\begin{Highlighting}[]
\NormalTok{viewof shor\_target = Inputs.select(SHOR\_TARGETS, \{}
\NormalTok{  label: "N to factor",}
\NormalTok{  format: t =\textgreater{} t.label.replace(/=.*/, "").trim()}
\NormalTok{\})}
\end{Highlighting}
\end{Shaded}

\begin{Shaded}
\begin{Highlighting}[]
\NormalTok{//| echo: false}
\NormalTok{shor\_computing = \{}
\NormalTok{  shor\_target}
\NormalTok{  const div = document.createElement(\textquotesingle{}div\textquotesingle{})}
\NormalTok{  div.className = \textquotesingle{}shor{-}computing\textquotesingle{}}
\NormalTok{  div.innerHTML = \textquotesingle{}\textless{}span class="shor{-}spinner"\textgreater{}\textless{}/span\textgreater{} running QPE circuit…\textquotesingle{}}
\NormalTok{  yield div}
\NormalTok{  await shor\_result}
\NormalTok{  return document.createElement(\textquotesingle{}span\textquotesingle{})}
\NormalTok{\}}
\end{Highlighting}
\end{Shaded}

\begin{Shaded}
\begin{Highlighting}[]
\NormalTok{//| echo: false}
\NormalTok{shor\_result = new Promise(resolve =\textgreater{} setTimeout(() =\textgreater{} \{}
\NormalTok{  const \{ N, a, aSq, n \} = shor\_target}
\NormalTok{  const t = 3}
\NormalTok{  const aPows = [a, aSq, 1]}
\NormalTok{  const workQ = Array.from(\{ length: n \}, (\_, i) =\textgreater{} t + i)}

\NormalTok{  let c = new ket.Circuit(t + n).x(t + n {-} 1)}
\NormalTok{  for (let i = 0; i \textless{} t; i++) c = c.h(i)}
\NormalTok{  for (let k = 0; k \textless{} t; k++) \{}
\NormalTok{    if (aPows[k] === 1) continue}
\NormalTok{    c = c.unitary(shor\_controlled(shor\_mulMatrix(aPows[k], N, n), n), k, ...workQ)}
\NormalTok{  \}}
\NormalTok{  // Inverse QFT on counting register}
\NormalTok{  for (let i = 0; i \textless{} Math.floor(t / 2); i++) c = c.swap(i, t {-} 1 {-} i)}
\NormalTok{  for (let j = 0; j \textless{} t; j++) \{}
\NormalTok{    for (let kk = j {-} 1; kk \textgreater{}= 0; kk{-}{-}) c = c.cu1({-}Math.PI / 2 ** (j {-} kk), kk, j)}
\NormalTok{    c = c.h(j)}
\NormalTok{  \}}

\NormalTok{  // Read counting register probabilities from statevector}
\NormalTok{  const sv = c.statevector()}
\NormalTok{  const countProbs = \{\}}
\NormalTok{  for (const [idx, amp] of sv) \{}
\NormalTok{    const p = amp.re * amp.re + amp.im * amp.im}
\NormalTok{    if (p \textless{} 1e{-}10) continue}
\NormalTok{    const bits = idx.toString(2).padStart(t + n, \textquotesingle{}0\textquotesingle{})}
\NormalTok{    const cb = bits.slice(0, t)}
\NormalTok{    countProbs[cb] = (countProbs[cb] || 0) + p}
\NormalTok{  \}}

\NormalTok{  // Period extraction via continued fractions}
\NormalTok{  const periods = new Set()}
\NormalTok{  for (const [bits, p] of Object.entries(countProbs)) \{}
\NormalTok{    if (p \textless{} 0.05) continue}
\NormalTok{    const m = parseInt(bits, 2)}
\NormalTok{    if (m === 0) continue}
\NormalTok{    for (let r = 1; r \textless{}= 2 * (1 \textless{}\textless{} t); r++)}
\NormalTok{      for (let s = 1; s \textless{} r; s++)}
\NormalTok{        if (Math.abs(s / r {-} m / (1 \textless{}\textless{} t)) \textless{} 0.5 / (1 \textless{}\textless{} t)) periods.add(r)}
\NormalTok{  \}}

\NormalTok{  // Factor via GCD}
\NormalTok{  const factors = []}
\NormalTok{  for (const r of [...periods].sort((x, y) =\textgreater{} x {-} y)) \{}
\NormalTok{    if (r \% 2 !== 0) continue}
\NormalTok{    const x = r === 4 ? aSq : r === 2 ? a : null}
\NormalTok{    if (!x || x === 1 || x === N {-} 1) continue}
\NormalTok{    const f1 = shor\_gcd(x {-} 1, N), f2 = shor\_gcd(x + 1, N)}
\NormalTok{    const fs = [f1, f2].filter(f =\textgreater{} f \textgreater{} 1 \&\& f \textless{} N)}
\NormalTok{    if (fs.length) factors.push(\{ r, p: fs[0], q: N / fs[0] \})}
\NormalTok{  \}}

\NormalTok{  resolve(\{ N, a, t, countProbs, factors \})}
\NormalTok{\}, 30))}
\end{Highlighting}
\end{Shaded}

\begin{Shaded}
\begin{Highlighting}[]
\NormalTok{//| echo: false}
\NormalTok{\{}
\NormalTok{  const \{ N, a, t, countProbs, factors \} = shor\_result}
\NormalTok{  const maxP = Math.max(...Object.values(countProbs))}
\NormalTok{  const lines = [\textasciigrave{}N = $\{N\},  base a = $\{a\},  counting bits t = $\{t\}\textasciigrave{}, \textasciigrave{}\textasciigrave{}, \textasciigrave{}QPE counting register:\textasciigrave{}]}
\NormalTok{  for (const [bits, p] of Object.entries(countProbs).sort((x,y) =\textgreater{} parseInt(x[0],2){-}parseInt(y[0],2))) \{}
\NormalTok{    const m = parseInt(bits, 2)}
\NormalTok{    const bar = \textquotesingle{}█\textquotesingle{}.repeat(Math.round(p / maxP * 24))}
\NormalTok{    lines.push(\textasciigrave{}  m=$\{String(m).padStart(2)\}  φ=$\{m\}/$\{1\textless{}\textless{}t\}  $\{(p*100).toFixed(0).padStart(3)\}\%  $\{bar\}\textasciigrave{})}
\NormalTok{  \}}
\NormalTok{  lines.push(\textasciigrave{}\textasciigrave{})}
\NormalTok{  if (factors.length) \{}
\NormalTok{    lines.push(\textasciigrave{}Period r = $\{factors[0].r\}\textasciigrave{})}
\NormalTok{    lines.push(\textasciigrave{}$\{N\} = $\{factors[0].p\} × $\{factors[0].q\}  ✓\textasciigrave{})}
\NormalTok{  \} else \{}
\NormalTok{    lines.push(\textasciigrave{}No factors found — try a different base\textasciigrave{})}
\NormalTok{  \}}
\NormalTok{  return pre(lines.join(\textquotesingle{}\textbackslash{}n\textquotesingle{}))}
\NormalTok{\}}
\end{Highlighting}
\end{Shaded}

The counting register peaks at multiples of \texttt{2\^{}t\ /\ r}. For
r=2 with t=3 precision bits: peaks at m=0 and m=4. For r=4: peaks at
m=0, 2, 4, 6. Continued fractions recovers r from any non-zero peak.

\section{What's next}\label{whats-next-6}

All the circuits so far assumed perfect gates. Chapter 9 covers noise:
what goes wrong on real hardware and how \texttt{ket} models it.

\part{Reality}

\chapter{Noise and Real Hardware}\label{sec-noise}

\begin{Shaded}
\begin{Highlighting}[]
\NormalTok{//| echo: false}
\NormalTok{ket = window.\_\_ket\_load}
\NormalTok{svg = (s) =\textgreater{} \{ const d = document.createElement("div"); d.innerHTML = s; return d; \}}
\NormalTok{pre = (s) =\textgreater{} \{ const d = document.createElement("pre"); d.textContent = s; return d; \}}
\end{Highlighting}
\end{Shaded}

Every circuit in the previous chapters ran perfectly. Real quantum
hardware doesn't.

Qubits are physical objects --- superconducting circuits, trapped ions,
photons. They interact with their environment, and that interaction
corrupts the quantum state. Understanding noise is the difference
between a textbook circuit and one that works in practice.

\section{Decoherence}\label{decoherence}

A qubit's coherence is its ability to maintain superposition. Two
timescales govern how quickly it's lost:

\textbf{\(T_1\) --- relaxation time}: the qubit spontaneously decays
from \(|1\rangle\) to \(|0\rangle\), like an excited atom emitting a
photon. Energy leaks to the environment.

\textbf{\(T_2\) --- dephasing time}: the phase relationship between
\(|0\rangle\) and \(|1\rangle\) randomises, collapsing superposition
without flipping the bit. Always \(T_2 \leq 2T_1\).

A circuit's total runtime must be well below both timescales or errors
accumulate faster than computation.

\section{Simulating noise}\label{simulating-noise}

Pass a noise profile to \texttt{run()} to apply \(T_1\)/\(T_2\) decay.
The same Bell state, clean vs noisy:

\begin{Shaded}
\begin{Highlighting}[]
\NormalTok{bell = new ket.Circuit(2).h(0).cx(0, 1)}
\NormalTok{svg(bell.simulate(\{ shots: 2000 \}).toSVG())}
\end{Highlighting}
\end{Shaded}

\begin{Shaded}
\begin{Highlighting}[]
\NormalTok{svg(bell.run(\{ shots: 2000, noise: \{ t1: 50, t2: 30 \} \}).toSVG())}
\end{Highlighting}
\end{Shaded}

The clean Bell state shows only \(|00\rangle\) and \(|11\rangle\). The
noisy run leaks into \(|01\rangle\) and \(|10\rangle\) --- decoherence
partially collapses the entanglement during the circuit.

Stronger decoherence (\(T_1\) and \(T_2\) measured in microseconds;
shorter means noisier):

\begin{Shaded}
\begin{Highlighting}[]
\NormalTok{svg(bell.run(\{ shots: 2000, noise: \{ t1: 10, t2: 5 \} \}).toSVG())}
\end{Highlighting}
\end{Shaded}

At \(T_1 = 10\,\mu\text{s}\), the correlation is almost entirely
destroyed. The histogram looks nearly uniform --- noise has randomised
the result.

\section{Circuit depth matters}\label{circuit-depth-matters}

A shallower circuit has less time to decohere. The same identity
operation at different depths:

\begin{Shaded}
\begin{Highlighting}[]
\NormalTok{shallow = new ket.Circuit(1).h(0).h(0)   // H·H = I, depth 2}
\NormalTok{svg(shallow.run(\{ shots: 2000, noise: \{ t1: 100, t2: 50 \} \}).toSVG())}
\end{Highlighting}
\end{Shaded}

\begin{Shaded}
\begin{Highlighting}[]
\NormalTok{deep = (() =\textgreater{} \{ let c = new ket.Circuit(1); for (let i = 0; i \textless{} 10; i++) c = c.h(0); return c \})()}
\NormalTok{svg(deep.run(\{ shots: 2000, noise: \{ t1: 100, t2: 50 \} \}).toSVG())}
\end{Highlighting}
\end{Shaded}

Both are mathematically identity --- they should always return
\(|0\rangle\). The deep circuit leaks into \(|1\rangle\) substantially
more. More time in the circuit means more exposure to noise.

\section{Real device profiles}\label{real-device-profiles}

\texttt{ket} ships with noise models calibrated from real hardware:

\begin{Shaded}
\begin{Highlighting}[]
\NormalTok{pre(Object.keys(ket.DEVICES).join(\textquotesingle{}\textbackslash{}n\textquotesingle{}))}
\end{Highlighting}
\end{Shaded}

\begin{Shaded}
\begin{Highlighting}[]
\NormalTok{svg(bell.run(\{ shots: 2000, noise: "harmony" \}).toSVG())}
\end{Highlighting}
\end{Shaded}

\begin{Shaded}
\begin{Highlighting}[]
\NormalTok{svg(bell.run(\{ shots: 2000, noise: "ibm\_sherbrooke" \}).toSVG())}
\end{Highlighting}
\end{Shaded}

Different hardware, different error profiles. IonQ's trapped-ion systems
tend to have higher gate fidelity but slower gate times. IBM's
superconducting systems are faster but noisier per gate. The device
string selects the full calibrated noise model including \(T_1\),
\(T_2\), gate error rates, and readout error.

\section{Readout error mitigation}\label{readout-error-mitigation}

Not all noise happens during the circuit. Measurement itself is
imperfect: a qubit in \(|0\rangle\) occasionally registers as 1, and
vice versa. This is \textbf{readout error} --- it's separate from
decoherence and can be corrected after the fact.

\texttt{run()} accepts \texttt{readoutMitigation:\ true} to apply a
simple matrix inversion correction:

\begin{Shaded}
\begin{Highlighting}[]
\NormalTok{svg(bell.run(\{ shots: 2000, noise: "harmony" \}).toSVG())}
\end{Highlighting}
\end{Shaded}

\begin{Shaded}
\begin{Highlighting}[]
\NormalTok{svg(bell.run(\{ shots: 2000, noise: "harmony", readoutMitigation: true \}).toSVG())}
\end{Highlighting}
\end{Shaded}

The mitigated histogram is closer to the ideal
\(|00\rangle\)/\(|11\rangle\) split. The correction is linear: the
simulator estimates the readout error matrix from the device profile and
inverts it. It doesn't fix decoherence --- only the measurement step ---
but on current hardware, readout error is often the dominant source of
error for short circuits.

\section{Error correction (the idea)}\label{error-correction-the-idea}

The solution to decoherence is \textbf{quantum error correction} (QEC):
encode one logical qubit across several physical qubits so that any
single error can be detected and corrected without measuring --- and
thus collapsing --- the data.

The simplest example: the three-qubit bit-flip code encodes
\(|0\rangle \to |000\rangle\) and \(|1\rangle \to |111\rangle\). A
single bit-flip is detectable by comparing pairs --- majority vote
restores the original.

Full QEC requires 50--1000 physical qubits per logical qubit depending
on the code and target error rate. The algorithms in Chapters 6--8
assume perfect logical qubits. Getting there is the work of the next
decade.

\section{Simulating at scale with
MPS}\label{simulating-at-scale-with-mps}

The chapters so far have used a statevector simulator: it stores all
\(2^n\) amplitudes in memory and applies gates as matrix
multiplications. For 30 qubits, that's \(2^{30} \approx 10^9\) complex
numbers --- a few gigabytes. For 50 qubits it's a petabyte. For 100
qubits, more memory than exists on Earth.

But many circuits of practical interest don't actually need all of that.
If a circuit has low entanglement, it has structure that can be
exploited. The \textbf{Matrix Product State} (MPS) representation stores
a chain of small tensors instead of one giant vector. The cost is
controlled by a single number: \(\chi\), the \textbf{bond dimension},
which equals the maximum Schmidt rank across any bipartition of the
qubits.

For a circuit with no entanglement, \(\chi = 1\). For a moderately
entangled circuit, \(\chi\) might be 10 or 100. Only for highly
entangled states does \(\chi\) grow to \(2^{n/2}\) --- and only then
does MPS hit the same wall as statevector.

\subsection{A 1024-qubit circuit}\label{a-1024-qubit-circuit}

A classical statevector simulator would need \(2^{1024}\) amplitudes to
represent the following circuit. That number has about 300 digits --- it
dwarfs the number of atoms in the observable universe. No computer will
ever store it.

Yet \texttt{runMps} returns in milliseconds:

\begin{Shaded}
\begin{Highlighting}[]
\NormalTok{result\_1024 = \{}
\NormalTok{  let c = new ket.Circuit(1024).h(0)}
\NormalTok{  for (let i = 0; i \textless{} 1023; i++) c = c.cx(i, i+1)}
\NormalTok{  return c.runMps(\{ shots: 200 \})}
\NormalTok{\}}
\end{Highlighting}
\end{Shaded}

\begin{Shaded}
\begin{Highlighting}[]
\NormalTok{pre(\textasciigrave{}backend: $\{result\_1024.backend\}   peakChi: $\{result\_1024.peakChi\}   qubits: $\{result\_1024.qubits\}\textasciigrave{})}
\end{Highlighting}
\end{Shaded}

\texttt{peakChi\ =\ 2}. That's why it's fast. The GHZ circuit --- H
followed by a chain of CNOTs --- produces a state with exactly two
Schmidt values across every bipartition. Despite spanning 1024 qubits,
the entanglement is as simple as a Bell state. MPS captures it in a
chain of 2×2 tensors.

The outcome is only ever \(|00\ldots0\rangle\) or \(|11\ldots1\rangle\).
Here's what the histogram looks like on a smaller GHZ (8 qubits) where
the bitstrings fit on screen:

\begin{Shaded}
\begin{Highlighting}[]
\NormalTok{ghz8\_circ = \{}
\NormalTok{  let c = new ket.Circuit(8).h(0)}
\NormalTok{  for (let i = 0; i \textless{} 7; i++) c = c.cx(i, i+1)}
\NormalTok{  return c}
\NormalTok{\}}
\NormalTok{svg(ghz8\_circ.simulate(\{ shots: 500 \}).toSVG())}
\end{Highlighting}
\end{Shaded}

\subsection{Entanglement entropy at
scale}\label{entanglement-entropy-at-scale}

The bond entropies make this precise:

\begin{Shaded}
\begin{Highlighting}[]
\NormalTok{ghz8\_circ.bondEntropies()}
\end{Highlighting}
\end{Shaded}

Every bond carries exactly 1 bit of entanglement entropy --- the maximum
for a two-outcome superposition. But \(\chi = 2\) is the minimum
non-trivial bond dimension. The entanglement is wide (spread across
every bond) but shallow (only two branches). MPS handles this
efficiently because the tensor dimensions stay small.

Compare a product state:

\begin{Shaded}
\begin{Highlighting}[]
\NormalTok{\{}
\NormalTok{  let prod = new ket.Circuit(8).h(0).h(1).h(2).h(3).h(4).h(5).h(6).h(7)}
\NormalTok{  return prod.bondEntropies()}
\NormalTok{\}}
\end{Highlighting}
\end{Shaded}

All zeros. No gate connects any qubit to any other, so no entanglement
crosses any bond. \(\chi = 1\) throughout.

\subsection{When MPS is fast --- and when it
isn't}\label{when-mps-is-fast-and-when-it-isnt}

The cost of MPS scales as \(O(n \chi^3)\) per gate layer. For
low-entanglement circuits, \(\chi\) stays small regardless of \(n\), and
the simulation is efficient. For high-entanglement circuits, \(\chi\)
grows and eventually the simulation becomes as hard as statevector.

\begin{Shaded}
\begin{Highlighting}[]
\NormalTok{// Low entanglement: GHZ on 50 qubits — chi stays at 2}
\NormalTok{mps\_ghz50 = \{}
\NormalTok{  let c = new ket.Circuit(50).h(0)}
\NormalTok{  for (let i = 0; i \textless{} 49; i++) c = c.cx(i, i+1)}
\NormalTok{  return c.runMps(\{ shots: 100 \})}
\NormalTok{\}}
\end{Highlighting}
\end{Shaded}

\begin{Shaded}
\begin{Highlighting}[]
\NormalTok{pre(\textasciigrave{}GHZ{-}50 — backend: $\{mps\_ghz50.backend\}   peakChi: $\{mps\_ghz50.peakChi\}\textasciigrave{})}
\end{Highlighting}
\end{Shaded}

\begin{Shaded}
\begin{Highlighting}[]
\NormalTok{// Higher entanglement: brickwork circuit on 20 qubits — chi grows}
\NormalTok{mps\_brick20 = \{}
\NormalTok{  let c = new ket.Circuit(20)}
\NormalTok{  for (let layer = 0; layer \textless{} 6; layer++) \{}
\NormalTok{    for (let i = 0; i \textless{} 19; i += 2) c = c.h(i).cx(i, i+1)}
\NormalTok{    for (let i = 1; i \textless{} 18; i += 2) c = c.h(i).cx(i, i+1)}
\NormalTok{  \}}
\NormalTok{  return c.runMps(\{ shots: 100, maxBond: 64 \})}
\NormalTok{\}}
\end{Highlighting}
\end{Shaded}

\begin{Shaded}
\begin{Highlighting}[]
\NormalTok{pre(\textasciigrave{}Brickwork{-}20 — backend: $\{mps\_brick20.backend\}   peakChi: $\{mps\_brick20.peakChi\}\textasciigrave{})}
\end{Highlighting}
\end{Shaded}

The GHZ chain on 50 qubits has \texttt{peakChi\ =\ 2} and runs
instantly. The brickwork circuit --- alternating layers of H and CX ---
builds up genuine entanglement and \texttt{peakChi} climbs. The
\texttt{maxBond} parameter caps the bond dimension and introduces an
approximation; without the cap, the simulation would be exact but
exponentially slow.

\subsection{MPS is not an approximation (for low-entanglement
circuits)}\label{mps-is-not-an-approximation-for-low-entanglement-circuits}

This is the key point: when \texttt{maxBond} is not set, \texttt{runMps}
is \textbf{exact}. For any circuit where \(\chi\) stays below a
manageable threshold --- GHZ states, many chemistry ansätze, QAOA at low
depth, any circuit with a brickwork structure that doesn't fully
scramble --- the MPS simulation gives the same result as a perfect
statevector simulation.

The exponential difficulty of simulating quantum systems is not
universal. It depends on the entanglement structure of the state.
Circuits with low entanglement are efficiently simulable classically ---
and that's not a bug, it's a feature: it tells you exactly how much
quantum advantage a given circuit has.

When \(\chi\) grows exponentially, you've found a circuit that is
genuinely hard to simulate classically. That's the regime where quantum
hardware offers real advantage --- and where MPS, too, runs out of room.

\section{What's next}\label{whats-next-7}

The final chapter applies quantum circuits to their original motivation:
simulating physics. Hamiltonians, Trotterisation, and the Variational
Quantum Eigensolver.

\chapter{Hamiltonians and Quantum Simulation}\label{sec-hamiltonian}

\begin{Shaded}
\begin{Highlighting}[]
\NormalTok{//| echo: false}
\NormalTok{ket = window.\_\_ket\_load}
\NormalTok{svg = (s) =\textgreater{} \{ const d = document.createElement("div"); d.innerHTML = s; return d; \}}
\NormalTok{pre = (s) =\textgreater{} \{ const d = document.createElement("pre"); d.textContent = s; return d; \}}
\end{Highlighting}
\end{Shaded}

Feynman's original motivation for quantum computers was simulation:
classical computers are exponentially slow at simulating quantum
systems. A quantum computer, being quantum itself, should be able to do
it efficiently.

A \textbf{Hamiltonian} is the operator that governs a quantum system's
energy and time evolution. Finding its lowest-energy state --- the
\emph{ground state} --- answers questions in chemistry, materials
science, and drug discovery that classical computers can't reach.

\section{Hamiltonians as Pauli
strings}\label{hamiltonians-as-pauli-strings}

Any Hamiltonian on \(n\) qubits can be written as a sum of Pauli
operator products:

\[H = \sum_i c_i \; P_i\]

where each \(P_i\) is a tensor product of \(I, X, Y, Z\) across qubits
--- a \textbf{Pauli string}. In \texttt{ket}, each term is
\texttt{\{\ ops:\ \textquotesingle{}ZZ\textquotesingle{},\ coeff:\ -1.0\ \}}
where the string encodes the Pauli on each qubit.

\section{Trotter simulation}\label{trotter-simulation}

To simulate time evolution
\(|\psi(t)\rangle = e^{-iHt}|\psi(0)\rangle\), we use
\textbf{Trotterisation}: split time into small steps \(\Delta t\) and
apply each Pauli term separately per step:

\[e^{-iHt} \approx \left(\prod_i e^{-ic_i P_i \Delta t}\right)^{N_\text{steps}}\]

Evolving a single qubit under \(H = X\) --- starting in \(|0\rangle\),
at \(t = \pi/2\) the state reaches \(|1\rangle\) with certainty (a full
population inversion):

\begin{Shaded}
\begin{Highlighting}[]
\NormalTok{H\_x = [\{ ops: \textquotesingle{}X\textquotesingle{}, coeff: 1.0 \}]}
\NormalTok{trot = ket.trotter(1, H\_x, Math.PI / 2, 8)}
\NormalTok{svg(trot.toSVG())}
\end{Highlighting}
\end{Shaded}

\begin{Shaded}
\begin{Highlighting}[]
\NormalTok{evolution = new ket.Circuit(1).compose(trot)}
\NormalTok{svg(evolution.simulate(\{ shots: 500 \}).toSVG())}
\end{Highlighting}
\end{Shaded}

\(|0\rangle\) has been driven to \(|1\rangle\) by the Hamiltonian ---
the quantum analogue of Rabi oscillations in a two-level atom.

\section{The Variational Quantum
Eigensolver}\label{the-variational-quantum-eigensolver}

\textbf{VQE} is a hybrid classical-quantum algorithm for ground state
energies, designed for near-term hardware with shallow, noise-tolerant
circuits.

Parameterise a circuit (the \emph{ansatz}), measure the expectation
value \(\langle\psi(\theta)|H|\psi(\theta)\rangle\), and use a classical
optimiser to minimise it. By the variational principle the minimum is an
upper bound on the true ground state energy.

For \(H = Z\) the ground state is \(|1\rangle\) with energy \(-1\):

\begin{Shaded}
\begin{Highlighting}[]
\NormalTok{H\_single = [\{ ops: \textquotesingle{}Z\textquotesingle{}, coeff: 1.0 \}]}
\NormalTok{ansatz\_1q = (params) =\textgreater{} new ket.Circuit(1).ry(params[0], 0)}
\NormalTok{result\_1q = ket.minimize(ansatz\_1q, H\_single, [0.1])}
\NormalTok{pre(\textasciigrave{}Ground state energy: $\{result\_1q.energy.toFixed(6)\}   (exact: {-}1.0)\textbackslash{}nOptimal \textbackslash{}u03b8:           $\{result\_1q.params[0].toFixed(6)\} rad  (= \textbackslash{}u03c0)\textasciigrave{})}
\end{Highlighting}
\end{Shaded}

The optimiser finds \(\theta = \pi\) --- \texttt{ry(π)} maps
\(|0\rangle \to |1\rangle\). Energy \(-1\), exact.

\section{A two-qubit molecule}\label{a-two-qubit-molecule}

A simplified model of the hydrogen molecule H₂ near its equilibrium bond
length fits in a two-qubit Hamiltonian. The coefficients come from
quantum chemistry:

\begin{Shaded}
\begin{Highlighting}[]
\NormalTok{H\_h2 = [}
\NormalTok{  \{ ops: \textquotesingle{}II\textquotesingle{}, coeff: {-}0.8105 \},}
\NormalTok{  \{ ops: \textquotesingle{}ZI\textquotesingle{}, coeff:  0.1722 \},}
\NormalTok{  \{ ops: \textquotesingle{}IZ\textquotesingle{}, coeff: {-}0.2228 \},}
\NormalTok{  \{ ops: \textquotesingle{}ZZ\textquotesingle{}, coeff:  0.1201 \},}
\NormalTok{  \{ ops: \textquotesingle{}XX\textquotesingle{}, coeff:  0.1686 \},}
\NormalTok{]}

\NormalTok{ansatz\_2q = (params) =\textgreater{}}
\NormalTok{  new ket.Circuit(2).ry(params[0], 0).ry(params[1], 1).cx(0, 1)}

\NormalTok{result\_2q = ket.minimize(ansatz\_2q, H\_h2, [0.1, 0.1])}
\NormalTok{pre(\textasciigrave{}H\textbackslash{}u2082 ground state energy: $\{result\_2q.energy.toFixed(6)\} Hartree\textbackslash{}nOptimal params:         [$\{result\_2q.params.map(p =\textgreater{} p.toFixed(3)).join(\textquotesingle{}, \textquotesingle{})\}]\textasciigrave{})}
\end{Highlighting}
\end{Shaded}

Ground state energy of H₂ from a quantum circuit, found by optimising
two rotation angles. For larger molecules --- where the classical
simulation becomes intractable --- this is the kind of calculation
quantum computers are expected to accelerate dramatically.

\section{QAOA --- optimization on
graphs}\label{qaoa-optimization-on-graphs}

The \textbf{Quantum Approximate Optimization Algorithm} applies the same
hybrid idea to combinatorial problems. Where VQE finds ground state
energies, QAOA finds approximate solutions to optimization problems like
Max-Cut: partition a graph's nodes into two sets to maximise the number
of edges between them.

The ansatz alternates between a problem Hamiltonian \(e^{-i\gamma H_C}\)
(which encodes the objective) and a mixing Hamiltonian
\(e^{-i\beta H_B}\) (which explores solutions). Depth \(p\) means \(p\)
rounds of each.

A 4-node square graph has two optimal cuts --- alternating nodes
\(0101\) and \(1010\), each severing all 4 edges:

\begin{Shaded}
\begin{Highlighting}[]
\NormalTok{qaoa\_edges = [[0,1],[1,2],[2,3],[3,0]]}
\NormalTok{H\_maxcut = ket.maxCutHamiltonian(4, qaoa\_edges)}
\end{Highlighting}
\end{Shaded}

\begin{Shaded}
\begin{Highlighting}[]
\NormalTok{// p=1 QAOA}
\NormalTok{ansatz\_qaoa1 = (params) =\textgreater{} ket.qaoa(4, qaoa\_edges, [params[0]], [params[1]])}
\NormalTok{result\_qaoa1 = ket.minimize(ansatz\_qaoa1, H\_maxcut, [0.3, 0.7])}
\NormalTok{svg(ansatz\_qaoa1(result\_qaoa1.params).simulate(\{ shots: 1000 \}).toSVG())}
\end{Highlighting}
\end{Shaded}

\begin{Shaded}
\begin{Highlighting}[]
\NormalTok{// p=2 QAOA — deeper circuit, better approximation}
\NormalTok{ansatz\_qaoa2 = (params) =\textgreater{} ket.qaoa(4, qaoa\_edges, [params[0], params[1]], [params[2], params[3]])}
\NormalTok{result\_qaoa2 = ket.minimize(ansatz\_qaoa2, H\_maxcut, [0.3, 0.6, 0.7, 0.4])}
\NormalTok{svg(ansatz\_qaoa2(result\_qaoa2.params).simulate(\{ shots: 1000 \}).toSVG())}
\end{Highlighting}
\end{Shaded}

At \(p=1\) the optimal states get a modest boost. At \(p=2\),
\texttt{0101} and \texttt{1010} dominate --- the circuit has learned to
concentrate probability on the two maximum cuts. More layers means a
better approximation; in the limit \(p \to \infty\), QAOA finds the
exact optimum.

\begin{Shaded}
\begin{Highlighting}[]
\NormalTok{pre(\textasciigrave{}p=1 energy: $\{result\_qaoa1.energy.toFixed(4)\}   p=2 energy: $\{result\_qaoa2.energy.toFixed(4)\}   (max cut = 4)\textasciigrave{})}
\end{Highlighting}
\end{Shaded}

The energy rises toward 4 with depth. This is the variational principle
at work on a combinatorial problem --- the same machinery as VQE,
applied to graphs instead of molecules.

\section{Where this is going}\label{where-this-is-going}

The circuits in this book run on a simulator. The algorithms are
correct; the hardware to run them at useful scale is still being built.

The near-term roadmap:

\begin{itemize}
\tightlist
\item
  \textbf{NISQ era} (now): 100--1000 noisy qubits. VQE, QAOA, quantum ML
  --- shallow circuits that tolerate noise.
\item
  \textbf{Early fault-tolerant} (\textasciitilde5--10 years): logical
  qubits via error correction. First practical advantage for quantum
  chemistry and optimisation.
\item
  \textbf{Full fault-tolerant} (\textasciitilde15+ years): millions of
  physical qubits. Shor's at cryptographic scale, full Hamiltonian
  simulation of useful molecules.
\end{itemize}

The mathematics in this book is complete. The engineering is the work of
the next decade.

\chapter{Backends and Classical Control}\label{sec-backends}

\begin{Shaded}
\begin{Highlighting}[]
\NormalTok{//| echo: false}
\NormalTok{ket = window.\_\_ket\_load}
\NormalTok{svg = (s) =\textgreater{} \{ const d = document.createElement("div"); d.innerHTML = s; return d; \}}
\NormalTok{pre = (s) =\textgreater{} \{ const d = document.createElement("pre"); d.textContent = s; return d; \}}
\end{Highlighting}
\end{Shaded}

The same \texttt{.simulate()} call handles circuits ranging from 2
qubits to 1000+. It does this by routing to different backends depending
on what the circuit contains:

\begin{longtable}[]{@{}
  >{\raggedright\arraybackslash}p{(\linewidth - 4\tabcolsep) * \real{0.3333}}
  >{\raggedright\arraybackslash}p{(\linewidth - 4\tabcolsep) * \real{0.3333}}
  >{\raggedright\arraybackslash}p{(\linewidth - 4\tabcolsep) * \real{0.3333}}@{}}
\toprule\noalign{}
\begin{minipage}[b]{\linewidth}\raggedright
Backend
\end{minipage} & \begin{minipage}[b]{\linewidth}\raggedright
When used
\end{minipage} & \begin{minipage}[b]{\linewidth}\raggedright
Cost
\end{minipage} \\
\midrule\noalign{}
\endhead
\bottomrule\noalign{}
\endlastfoot
\textbf{Clifford} & Circuits using only H, S, X, Y, Z, CNOT, CZ, SWAP &
\(O(n^2)\) per gate --- handles 1000+ qubits \\
\textbf{Statevector} & General circuits, no mid-circuit measurement &
\(O(2^n)\) memory --- exact, small \(n\) \\
\textbf{MPS} & Low-entanglement circuits via \texttt{runMps()} &
\(O(n \chi^3)\) --- scales with entanglement \\
\textbf{Density matrix} & Any circuit via \texttt{run(\{\ noise\ \})} &
\(O(4^n)\) --- exact noise simulation \\
\end{longtable}

The routing is automatic. You don't select a backend ---
\texttt{simulate()} inspects the circuit and picks the right one.

\section{Clifford circuits at scale}\label{clifford-circuits-at-scale}

A \textbf{Clifford circuit} uses only stabilizer gates: H, S, X, Y, Z,
CNOT, CZ, SWAP, and measurement. The Gottesman--Knill theorem says these
circuits are classically simulable in polynomial time using a
\emph{stabilizer tableau} --- a compact \(O(n^2)\) representation that
never touches the exponential state space.

\begin{Shaded}
\begin{Highlighting}[]
\NormalTok{// 1000{-}qubit circuit — every qubit put in superposition}
\NormalTok{clifford\_1000 = \{}
\NormalTok{  let c = new ket.Circuit(1000)}
\NormalTok{  for (let i = 0; i \textless{} 1000; i++) c = c.h(i)}
\NormalTok{  return c.simulate(\{ shots: 200 \})}
\NormalTok{\}}
\end{Highlighting}
\end{Shaded}

\begin{Shaded}
\begin{Highlighting}[]
\NormalTok{pre(\textasciigrave{}backend: $\{clifford\_1000.backend\}   qubits: $\{clifford\_1000.qubits\}\textasciigrave{})}
\end{Highlighting}
\end{Shaded}

A 1000-qubit statevector would need \(2^{1000}\) amplitudes --- a number
with 300 digits. The Clifford backend handles it by never constructing
that vector at all.

The price: Clifford circuits can't produce arbitrary quantum states. The
Clifford group is exactly the set of circuits that \emph{don't} require
exponential classical resources to simulate. This is why random Clifford
circuits don't give quantum advantage --- they sit inside the simulable
region. Quantum speedup requires gates that escape it (like T, Toffoli,
and the gates inside Shor's controlled-\(U_a\)).

Here's the same 1000-qubit structure but with a CNOT chain, which is
still Clifford:

\begin{Shaded}
\begin{Highlighting}[]
\NormalTok{clifford\_ghz = \{}
\NormalTok{  let c = new ket.Circuit(1000).h(0)}
\NormalTok{  for (let i = 0; i \textless{} 999; i++) c = c.cx(i, i+1)}
\NormalTok{  return c.simulate(\{ shots: 200 \})}
\NormalTok{\}}
\end{Highlighting}
\end{Shaded}

\begin{Shaded}
\begin{Highlighting}[]
\NormalTok{pre(\textasciigrave{}backend: $\{clifford\_ghz.backend\}   qubits: $\{clifford\_ghz.qubits\}\textasciigrave{})}
\end{Highlighting}
\end{Shaded}

Only \(|00\ldots0\rangle\) and \(|11\ldots1\rangle\) --- a 1000-qubit
GHZ state, simulated classically in seconds.

\section{Mid-circuit measurement and classical
control}\label{mid-circuit-measurement-and-classical-control}

Every circuit so far measured at the end. \texttt{ket} supports
measuring qubits \textbf{mid-circuit} and using the classical result to
conditionally apply gates --- the same feedback loop used in quantum
error correction and teleportation.

\begin{Shaded}
\begin{Highlighting}[]
\NormalTok{// Measure |+⟩, then conditionally flip a second qubit to match}
\NormalTok{c\_midmeas = new ket.Circuit(2)}
\NormalTok{  .h(0)}
\NormalTok{  .measure(0, \textquotesingle{}m\textquotesingle{}, 0)}
\NormalTok{  .if(\textquotesingle{}m\textquotesingle{}, 1, c =\textgreater{} c.x(1))}
\NormalTok{svg(c\_midmeas.toSVG())}
\end{Highlighting}
\end{Shaded}

\begin{Shaded}
\begin{Highlighting}[]
\NormalTok{svg(c\_midmeas.simulate(\{ shots: 1000 \}).toSVG())}
\end{Highlighting}
\end{Shaded}

Qubit 0 is measured. If it came out 1, qubit 1 is flipped. The result is
always \(|00\rangle\) or \(|11\rangle\) --- the two qubits are forced to
agree by classical feedback, even though they were never entangled.

\subsection{Quantum teleportation}\label{quantum-teleportation}

The canonical application of mid-circuit measurement is \textbf{quantum
teleportation}: transmitting a qubit's full quantum state using only
classical bits and a pre-shared entangled pair.

The protocol:

\begin{enumerate}
\def\labelenumi{\arabic{enumi}.}
\tightlist
\item
  Alice and Bob share a Bell pair (qubits 1 and 2)
\item
  Alice holds the qubit to teleport (qubit 0)
\item
  Alice entangles her qubit with her half of the Bell pair and measures
  both
\item
  Alice sends her two classical measurement bits to Bob
\item
  Bob applies conditional X and Z corrections to his qubit
\end{enumerate}

The state of qubit 0 arrives at qubit 2 --- no quantum channel, just two
classical bits.

\begin{Shaded}
\begin{Highlighting}[]
\NormalTok{// Teleport |1⟩ from qubit 0 to qubit 2}
\NormalTok{tele = new ket.Circuit(3)}
\NormalTok{  .x(0)                         // state to teleport: |1⟩}
\NormalTok{  .h(1).cx(1, 2)                // Bell pair on qubits 1 and 2}
\NormalTok{  .cx(0, 1).h(0)                // Alice\textquotesingle{}s Bell{-}basis measurement}
\NormalTok{  .measure(0, \textquotesingle{}m0\textquotesingle{}, 0)}
\NormalTok{  .measure(1, \textquotesingle{}m1\textquotesingle{}, 0)}
\NormalTok{  .if(\textquotesingle{}m1\textquotesingle{}, 1, c =\textgreater{} c.x(2))    // Bob\textquotesingle{}s correction: X if m1 = 1}
\NormalTok{  .if(\textquotesingle{}m0\textquotesingle{}, 1, c =\textgreater{} c.z(2))    // Bob\textquotesingle{}s correction: Z if m0 = 1}
\NormalTok{svg(tele.toSVG())}
\end{Highlighting}
\end{Shaded}

\begin{Shaded}
\begin{Highlighting}[]
\NormalTok{svg(tele.simulate(\{ shots: 1000 \}).toSVG())}
\end{Highlighting}
\end{Shaded}

Alice's measurement outcomes (\(m_0\), \(m_1\)) are random and uniformly
distributed across all four combinations. But qubit 2 --- Bob's qubit
--- is always \(|1\rangle\) regardless of which outcome Alice got. The
corrections steer it to the right state every time.

This is what makes teleportation remarkable: the quantum information
travels without a quantum channel. The two classical bits tell Bob which
of four correction operations to apply, and after applying it, he holds
the original state exactly.

\section{Quantum error correction}\label{quantum-error-correction}

Mid-circuit measurement makes \textbf{quantum error correction} (QEC)
possible. The idea: encode one logical qubit across several physical
qubits, measure \emph{parities} (not values) to detect errors, then
correct without ever touching the logical state.

The simplest code is the \textbf{3-qubit bit-flip code}:

\begin{itemize}
\tightlist
\item
  Encode: \(|\psi\rangle = \alpha|0\rangle + \beta|1\rangle\) →
  \(\alpha|000\rangle + \beta|111\rangle\)
\item
  Two syndrome ancillas measure \(Z_0Z_1\) and \(Z_1Z_2\) --- the parity
  of adjacent qubit pairs
\item
  Any single bit-flip changes exactly two parities, uniquely identifying
  which qubit flipped
\item
  Classical control corrects the right qubit; the logical state is
  unchanged
\end{itemize}

\begin{Shaded}
\begin{Highlighting}[]
\NormalTok{// Qubits 0,1,2 = logical data  |  qubits 3,4 = syndrome ancillas}
\NormalTok{qec\_error = new ket.Circuit(5)}
\NormalTok{  .h(0)                                      // logical |+⟩ to protect}
\NormalTok{  .cx(0, 1).cx(0, 2)                         // encode}
\NormalTok{  .x(1)                                      // inject error: flip qubit 1}
\NormalTok{  .cx(0, 3).cx(1, 3).measure(3, \textquotesingle{}s\textquotesingle{}, 0)     // syndrome: parity(q0,q1) → s[0]}
\NormalTok{  .cx(1, 4).cx(2, 4).measure(4, \textquotesingle{}s\textquotesingle{}, 1)     // syndrome: parity(q1,q2) → s[1]}
\NormalTok{  .if(\textquotesingle{}s\textquotesingle{}, 1, c =\textgreater{} c.x(0))                  // s=01 → correct q0}
\NormalTok{  .if(\textquotesingle{}s\textquotesingle{}, 3, c =\textgreater{} c.x(1))                  // s=11 → correct q1}
\NormalTok{  .if(\textquotesingle{}s\textquotesingle{}, 2, c =\textgreater{} c.x(2))                  // s=10 → correct q2}
\NormalTok{  .cx(0, 1).cx(0, 2)                        // decode}
\NormalTok{svg(qec\_error.toSVG())}
\end{Highlighting}
\end{Shaded}

\begin{Shaded}
\begin{Highlighting}[]
\NormalTok{svg(qec\_error.simulate(\{ shots: 1000 \}).toSVG())}
\end{Highlighting}
\end{Shaded}

The histogram shows only two outcomes --- qubit 0 is 50/50 (the original
\(|+\rangle\) recovered), qubits 1 and 2 are always \(|0\rangle\) after
decoding, and the syndrome ancillas (qubits 3 and 4) are always
\(|11\rangle\), confirming that a qubit-1 error was detected and
corrected every shot.

Compare the clean case --- no error injected:

\begin{Shaded}
\begin{Highlighting}[]
\NormalTok{qec\_clean = new ket.Circuit(5)}
\NormalTok{  .h(0)}
\NormalTok{  .cx(0, 1).cx(0, 2)}
\NormalTok{  .cx(0, 3).cx(1, 3).measure(3, \textquotesingle{}s\textquotesingle{}, 0)}
\NormalTok{  .cx(1, 4).cx(2, 4).measure(4, \textquotesingle{}s\textquotesingle{}, 1)}
\NormalTok{  .if(\textquotesingle{}s\textquotesingle{}, 1, c =\textgreater{} c.x(0))}
\NormalTok{  .if(\textquotesingle{}s\textquotesingle{}, 3, c =\textgreater{} c.x(1))}
\NormalTok{  .if(\textquotesingle{}s\textquotesingle{}, 2, c =\textgreater{} c.x(2))}
\NormalTok{  .cx(0, 1).cx(0, 2)}
\NormalTok{svg(qec\_clean.simulate(\{ shots: 1000 \}).toSVG())}
\end{Highlighting}
\end{Shaded}

Syndrome ancillas are \(|00\rangle\) --- no error detected. Logical
qubit 0 is still \(|+\rangle\). The parity measurements extracted error
information without disturbing the encoded data.

This is the principle behind every fault-tolerant quantum computer:
errors are detectable and correctable because the logical state is
spread across multiple qubits, and syndromes can be measured without
collapsing the logical superposition.

\section{Custom noise channels}\label{custom-noise-channels}

The T1/T2 model from Chapter 9 describes thermal relaxation and
dephasing. But real hardware has many other error types. \texttt{run()}
accepts custom \textbf{Kraus operators} --- the most general way to
describe any quantum channel.

A Kraus channel \(\mathcal{E}(\rho) = \sum_k K_k \rho K_k^\dagger\) is
valid whenever \(\sum_k K_k^\dagger K_k = I\). The simulator applies the
channel after each single-qubit gate (pass \texttt{kraus1}) or two-qubit
gate (pass \texttt{kraus2}).

\textbf{Amplitude damping} models energy loss to the environment --- a
qubit in \(|1\rangle\) spontaneously decays to \(|0\rangle\) with
probability \(\gamma\):

\[K_0 = \begin{pmatrix} 1 & 0 \\ 0 & \sqrt{1-\gamma} \end{pmatrix} \qquad K_1 = \begin{pmatrix} 0 & \sqrt{\gamma} \\ 0 & 0 \end{pmatrix}\]

\begin{Shaded}
\begin{Highlighting}[]
\NormalTok{c\_kraus = (re, im = 0) =\textgreater{} (\{ re, im \})}

\NormalTok{// Amplitude damping: qubit decays from |1⟩ to |0⟩ with probability gamma per gate}
\NormalTok{gamma\_low  = 0.05}
\NormalTok{gamma\_high = 0.40}

\NormalTok{K0\_low  = [[c\_kraus(1), c\_kraus(0)], [c\_kraus(0), c\_kraus(Math.sqrt(1 {-} gamma\_low))]]}
\NormalTok{K1\_low  = [[c\_kraus(0), c\_kraus(Math.sqrt(gamma\_low))],  [c\_kraus(0), c\_kraus(0)]]}
\NormalTok{K0\_high = [[c\_kraus(1), c\_kraus(0)], [c\_kraus(0), c\_kraus(Math.sqrt(1 {-} gamma\_high))]]}
\NormalTok{K1\_high = [[c\_kraus(0), c\_kraus(Math.sqrt(gamma\_high))], [c\_kraus(0), c\_kraus(0)]]}
\end{Highlighting}
\end{Shaded}

\begin{Shaded}
\begin{Highlighting}[]
\NormalTok{// Start in |1⟩ — amplitude damping pulls it toward |0⟩}
\NormalTok{circ\_damp = new ket.Circuit(1).x(0)}
\NormalTok{svg(circ\_damp.run(\{ shots: 2000, noise: \{ kraus1: [K0\_low,  K1\_low]  \} \}).toSVG())}
\end{Highlighting}
\end{Shaded}

\begin{Shaded}
\begin{Highlighting}[]
\NormalTok{svg(circ\_damp.run(\{ shots: 2000, noise: \{ kraus1: [K0\_high, K1\_high] \} \}).toSVG())}
\end{Highlighting}
\end{Shaded}

At \(\gamma = 0.05\) most population stays in \(|1\rangle\). At
\(\gamma = 0.40\) the decay is severe --- the energy loss dominates.
Amplitude damping is not captured by dephasing alone: T2 noise
randomises the phase but doesn't move population between \(|0\rangle\)
and \(|1\rangle\). Kraus operators let you model both independently.

\textbf{Phase damping} --- pure dephasing with no energy loss --- has
its own Kraus form. You can also describe coherent errors, leakage to
non-computational states, or any other physically motivated noise. The
constraint is just \(\sum_k K_k^\dagger K_k = I\).

\section{What this means}\label{what-this-means}

The backend architecture reflects a real partition in quantum computing:

\begin{itemize}
\tightlist
\item
  \textbf{Clifford circuits} are classically simulable. They don't give
  quantum advantage. They're useful for error correction codes and
  benchmarking --- but not for computation.
\item
  \textbf{Non-Clifford circuits} (T gates, Toffoli, controlled
  rotations) escape classical simulation and are where quantum advantage
  lives.
\item
  \textbf{Low-entanglement circuits} (small \(\chi\)) can be handled by
  MPS regardless of qubit count.
\item
  \textbf{Noisy circuits} require density matrix simulation --- the
  state is now a \(2^n \times 2^n\) matrix, not a vector.
\end{itemize}

Mid-circuit measurement and classical control bridge classical and
quantum computation --- they're how error-corrected logical qubits are
implemented in practice, where syndrome measurements drive real-time
correction without disturbing the data.

\bookmarksetup{startatroot}

\chapter*{API Reference}\label{sec-appendix}
\addcontentsline{toc}{chapter}{API Reference}

\markboth{API Reference}{API Reference}

\section*{Circuit --- gates}\label{circuit-gates}
\addcontentsline{toc}{section}{Circuit --- gates}

\markright{Circuit --- gates}

\begin{longtable}[]{@{}
  >{\raggedright\arraybackslash}p{(\linewidth - 2\tabcolsep) * \real{0.5000}}
  >{\raggedright\arraybackslash}p{(\linewidth - 2\tabcolsep) * \real{0.5000}}@{}}
\toprule\noalign{}
\begin{minipage}[b]{\linewidth}\raggedright
Method
\end{minipage} & \begin{minipage}[b]{\linewidth}\raggedright
Description
\end{minipage} \\
\midrule\noalign{}
\endhead
\bottomrule\noalign{}
\endlastfoot
\texttt{new\ Circuit(n)} & Create an n-qubit circuit \\
\texttt{.h(q)} & Hadamard \\
\texttt{.x(q)} / \texttt{.y(q)} / \texttt{.z(q)} & Pauli gates \\
\texttt{.s(q)} / \texttt{.t(q)} & Phase / T gate \\
\texttt{.rx(θ,\ q)} / \texttt{.ry(θ,\ q)} / \texttt{.rz(θ,\ q)} &
Rotation gates --- angle first \\
\texttt{.cx(ctrl,\ tgt)} & CNOT \\
\texttt{.cz(ctrl,\ tgt)} / \texttt{.cs(ctrl,\ tgt)} & Controlled-Z /
Controlled-S \\
\texttt{.ccx(c1,\ c2,\ tgt)} & Toffoli (doubly controlled X) \\
\texttt{.swap(a,\ b)} & SWAP \\
\texttt{.measure(q?)} & Measure qubit(s); omit for all \\
\texttt{.compose(other)} & Append another circuit \\
\end{longtable}

\section*{Circuit --- simulation}\label{circuit-simulation}
\addcontentsline{toc}{section}{Circuit --- simulation}

\markright{Circuit --- simulation}

\begin{longtable}[]{@{}
  >{\raggedright\arraybackslash}p{(\linewidth - 2\tabcolsep) * \real{0.5000}}
  >{\raggedright\arraybackslash}p{(\linewidth - 2\tabcolsep) * \real{0.5000}}@{}}
\toprule\noalign{}
\begin{minipage}[b]{\linewidth}\raggedright
Method
\end{minipage} & \begin{minipage}[b]{\linewidth}\raggedright
Returns
\end{minipage} \\
\midrule\noalign{}
\endhead
\bottomrule\noalign{}
\endlastfoot
\texttt{.simulate(\{\ shots\ \})} & \texttt{Distribution} ---
statevector/Clifford auto-routing \\
\texttt{.run(\{\ shots,\ noise,\ readoutMitigation?\ \})} &
\texttt{Distribution} --- density matrix with noise \\
\texttt{.runMps(\{\ shots,\ maxBond?\ \})} &
\texttt{\{\ backend,\ peakChi,\ probs,\ qubits\ \}} --- \texttt{probs}
is a plain \texttt{\{\ bitstring:\ prob\ \}} map \\
\texttt{.bondEntropies()} & \texttt{number{[}{]}} --- log₂ von Neumann
entropy per bond \\
\texttt{.stateAsArray()} &
\texttt{{[}\{\ bitstring,\ re,\ im,\ prob,\ phase\ \}{]}} \\
\texttt{.toSVG()} & Circuit diagram as SVG string \\
\texttt{.draw()} & Circuit diagram as ASCII string \\
\texttt{.blochSphere(q)} & Bloch sphere SVG for qubit q \\
\texttt{.blochAngles(q)} & \texttt{\{\ theta,\ phi\ \}} in radians \\
\end{longtable}

\section*{Distribution}\label{distribution}
\addcontentsline{toc}{section}{Distribution}

\markright{Distribution}

\begin{longtable}[]{@{}ll@{}}
\toprule\noalign{}
Property / Method & Description \\
\midrule\noalign{}
\endhead
\bottomrule\noalign{}
\endlastfoot
\texttt{.probs} & \texttt{\{\ bitstring:\ probability\ \}} object \\
\texttt{.toSVG()} & Histogram as SVG string \\
\texttt{.backend} & Which backend ran the circuit \\
\end{longtable}

\section*{Algorithms}\label{algorithms-1}
\addcontentsline{toc}{section}{Algorithms}

\markright{Algorithms}

\begin{longtable}[]{@{}
  >{\raggedright\arraybackslash}p{(\linewidth - 4\tabcolsep) * \real{0.3333}}
  >{\raggedright\arraybackslash}p{(\linewidth - 4\tabcolsep) * \real{0.3333}}
  >{\raggedright\arraybackslash}p{(\linewidth - 4\tabcolsep) * \real{0.3333}}@{}}
\toprule\noalign{}
\begin{minipage}[b]{\linewidth}\raggedright
Symbol
\end{minipage} & \begin{minipage}[b]{\linewidth}\raggedright
Signature
\end{minipage} & \begin{minipage}[b]{\linewidth}\raggedright
Description
\end{minipage} \\
\midrule\noalign{}
\endhead
\bottomrule\noalign{}
\endlastfoot
\texttt{grover} & \texttt{(n,\ oracle,\ iters?)} & Grover's search;
oracle is \texttt{(c)\ =\textgreater{}\ c} \\
\texttt{trotter} & \texttt{(n,\ H,\ t,\ steps)} & Trotterised
Hamiltonian simulation; H is \texttt{{[}\{ops,\ coeff\}{]}} \\
\texttt{minimize} & \texttt{(ansatz,\ H,\ params)} & VQE; ansatz is
\texttt{(params)\ =\textgreater{}\ Circuit} \\
\texttt{shorBeauregard} & \texttt{(N,\ a)} (BigInt) & Beauregard Shor's;
returns \texttt{\{\ factors,\ a,\ qubits\ \}} \\
\texttt{factor} & \texttt{(N)} (BigInt) & Convenience wrapper around
Shor's \\
\texttt{qft} & \texttt{(circuit,\ qubits{[}{]})} & In-place QFT on named
qubits \\
\texttt{applyQft} & \texttt{(circuit,\ n,\ offset)} & QFT on n qubits at
offset \\
\texttt{applyIqft} & \texttt{(circuit,\ n,\ offset)} & Inverse QFT on n
qubits at offset \\
\texttt{DEVICES} & --- & Object of calibrated hardware noise profiles \\
\end{longtable}

\section*{Math notation glossary}\label{math-notation-glossary}
\addcontentsline{toc}{section}{Math notation glossary}

\markright{Math notation glossary}

\begin{longtable}[]{@{}
  >{\raggedright\arraybackslash}p{(\linewidth - 2\tabcolsep) * \real{0.5000}}
  >{\raggedright\arraybackslash}p{(\linewidth - 2\tabcolsep) * \real{0.5000}}@{}}
\toprule\noalign{}
\begin{minipage}[b]{\linewidth}\raggedright
Symbol
\end{minipage} & \begin{minipage}[b]{\linewidth}\raggedright
Meaning
\end{minipage} \\
\midrule\noalign{}
\endhead
\bottomrule\noalign{}
\endlastfoot
\(\vert\psi\rangle\) & Quantum state (ket) \\
\(\langle\psi\vert\) & Conjugate transpose (bra) \\
\(\alpha, \beta \in \mathbb{C}\) & Complex amplitudes \\
\(\vert\alpha\vert^2\) & Measurement probability \\
\(\otimes\) & Tensor product \\
\(U^\dagger\) & Conjugate transpose (inverse unitary) \\
\(\langle H\rangle = \langle\psi\vert H\vert\psi\rangle\) & Expectation
value \\
\(\chi\) & MPS bond dimension \\
\(S_\text{vN}\) & Von Neumann entanglement entropy \\
\end{longtable}


\backmatter


\end{document}
